% concatenated from morley_tetrahedron.tex
\documentclass[11pt]{article}

\usepackage{amsmath,amssymb,amsthm}
\usepackage[a4paper,margin=1in]{geometry}
\usepackage{booktabs}
\usepackage{enumitem}
\usepackage{microtype}
\usepackage[title,titletoc]{appendix}
\usepackage[hidelinks]{hyperref}

\newtheorem{theorem}{Theorem}[section]
\newtheorem{proposition}[theorem]{Proposition}
\newtheorem{lemma}[theorem]{Lemma}
\newtheorem{corollary}[theorem]{Corollary}
\newtheorem{conjecture}[theorem]{Conjecture}
\theoremstyle{definition}
\newtheorem{example}[theorem]{Example}
\theoremstyle{remark}
\newtheorem{remark}[theorem]{Remark}
\newtheorem{problem}[theorem]{Problem}

\DeclareMathOperator{\tr}{tr}
\DeclareMathOperator{\spn}{span}
\DeclareMathOperator{\dist}{dist}
\DeclareMathOperator{\vol}{vol}
\DeclareMathOperator{\diag}{diag}
\DeclareMathOperator{\rank}{rank}
\DeclareMathOperator{\Res}{Res}
\DeclareMathOperator{\pp}{pp}
\newcommand{\HS}{\mathrm{HS}}
\newcommand{\T}{\mathsf{T}}
\newcommand{\R}{\mathbb{R}}
\newcommand{\bu}{\mathbf u}
\newcommand{\be}{\mathbf e}
\newcommand{\bv}{\mathbf v}
\newcommand{\bw}{\mathbf w}
\newcommand{\bz}{\mathbf z}
\newcommand{\bq}{\mathbf q}
\newcommand{\bff}{\mathbf f}
\newcommand{\bzero}{\mathbf 0}
\newcommand{\norm}[1]{\lVert #1 \rVert}
\newcommand{\M}{\mathcal M}
\newcommand{\one}{\mathbf 1}
\newcommand{\ph}{\varphi}
\allowdisplaybreaks[1]
\raggedbottom
\renewcommand{\labelitemi}{$\bullet$}
\setcounter{tocdepth}{1}
\hypersetup{pdftitle={The Converse Problem for the Morley Tetrahedron},pdfauthor={Quang Hung Tran}}
\renewcommand{\contentsname}{Contents}

\title{The Converse Problem for the Morley Tetrahedron:\\Counterexamples, Conjectures, and Partial Results}
\author{Quang Hung Tran}
\date{}

\begin{document}
\maketitle

\begin{abstract}
An earlier paper in Acta Mathematica Hungarica showed that the Morley construction preserves equality of opposite edge pairs and proposed two converse conjectures. We disprove both. A nonisosceles tetrahedron $T_1$ and an isosceles but nonregular tetrahedron $T_2$ have regular Morley tetrahedra, and an analytic curve of nonisosceles tetrahedra has isosceles Morley tetrahedra. The tetrahedron $T_2$ has edges $AB=CD=1$ and $AC=AD=BC=BD=\sqrt{(21+4\sqrt6)/45}$. Both examples have four equal cross edges. We conjecture that every tetrahedron with a regular Morley tetrahedron has this property, and prove it whenever the original tetrahedron has a nontrivial symmetry. Within the class with four equal cross edges, only the regular tetrahedron, $T_1$ and $T_2$ have regular Morley tetrahedra. The sextic defining $T_1$ has Galois group $S_6$, so $T_1$ cannot be expressed by radicals. The proof-critical computer-assisted checks use exact rational arithmetic.
\end{abstract}

\noindent
\textbf{Keywords.} Morley's theorem, Morley tetrahedron, isosceles tetrahedron, disphenoid, dihedral angle, reflection symmetry, interval arithmetic.

\medskip
\noindent
\textbf{2020 Mathematics Subject Classification.}
Primary 51M04; Secondary 51M20, 12F10, 65G30.

\tableofcontents

\section{Introduction}

Morley's trisector theorem states that the intersections of adjacent angle trisectors of a triangle form an equilateral triangle. The theorem goes back to Morley \cite{Morley}; see also Child \cite{Child}, Coxeter and Greitzer \cite{CoxeterGreitzer}, Oakley and Baker \cite{OakleyBaker}, Hashimoto \cite{Hashimoto}, and Grinberg and Orhon \cite{GrinbergOrhon} for various proofs and viewpoints. Several extensions of Morley's configuration have also been investigated; see, for example, \cite{Braude,DergiadesTran}. A recent synthetic proof is given in \cite{TranSynthetic}.

A three dimensional analogue was considered by Svrtan and Veljan \cite{SvrtanVeljan}. In \cite{TranActa} the author defined a Morley tetrahedron by trisecting the six dihedral angles of a tetrahedron.

We call a tetrahedron $ABCD$ \emph{isosceles} when its three pairs of opposite edges have equal lengths,
\[
AB=CD,\; AC=BD,\; AD=BC;
\]
regular tetrahedra are included. Definition~1 of \cite{TranActa} excluded regular tetrahedra from this terminology, but the proof of its direct theorem establishes the equalities of opposite edge pairs. We use that edge-equality statement here. Such tetrahedra are also called \emph{equifacial tetrahedra} or \emph{disphenoids}; see \cite{BrandtsKrizek,HajjaWalker,Leech}. Related tetrahedra with congruent pairs of facets are studied in \cite{Klain,TranVolume}.

\medskip

\noindent
\textbf{The direct theorem, in edge-equality form \cite{TranActa}.} If $ABCD$ is isosceles, then its Morley tetrahedron is isosceles.

\medskip

In \cite{TranActa} the author also proposed two converse statements.

\begin{conjecture}[The original converse conjecture]
\label{conj:old1} If the Morley tetrahedron of $ABCD$ is regular, respectively isosceles, then $ABCD$ is regular, respectively isosceles.
\end{conjecture}

Theorem 6 of \cite{TranActa} gives a more general construction, in which the trisecting planes are replaced by suitable $p$ secting planes depending on parameters $e,f,g>1$. The corresponding converse was proposed as well.

\begin{conjecture}[The original $p$ sector converse]
\label{conj:old2} The converse statement analogous to Conjecture~\ref{conj:old1} holds for the generalized $p$ sector tetrahedron of \cite[Theorem 6]{TranActa}.
\end{conjecture}

\subsection*{Results}

The paper is divided into three parts.

Part~\ref{part:counterexamples} contains the counterexamples. Theorem~\ref{thm:T1} gives a nonisosceles tetrahedron $T_1$ with a regular Morley tetrahedron, and Theorem~\ref{thm:T2} gives an isosceles one, $T_2$, in closed form. Theorem~\ref{thm:curve} gives a curve of nonisosceles tetrahedra whose Morley tetrahedra are isosceles but not regular. The generalized $p$ sector construction of \cite{TranActa} includes ordinary trisection at $e=f=g=3$, so both converse conjectures fail.

In Part~\ref{part:conjectures} we look for a statement that could replace them. Both $T_1$ and $T_2$ have a pair of opposite edges, say $AB$ and $CD$, such that the four cross edges are equal. We conjecture that every tetrahedron with a regular Morley tetrahedron has this property (Conjecture~\ref{conj:four}), and we discuss several variants.

Part~\ref{part:results} collects what we can prove. Suppose that the Morley tetrahedron of $ABCD$ is regular. If $ABCD$ is isosceles, or has two pairs of equal opposite edges, or has a reflection symmetry, then it has four equal cross edges (Theorems~\ref{thm:disphenoid}, \ref{thm:twopairs} and~\ref{thm:reflection}). If it has four equal cross edges, it is similar to the regular tetrahedron, to $T_1$ or to $T_2$ (Theorem~\ref{sym:thm:classification}). Together these results show that a tetrahedron with a nontrivial symmetry and a regular Morley tetrahedron is one of these three (Corollary~\ref{cor:symmetric}). A counterexample to Conjecture~\ref{conj:four}, if there is one, therefore has no symmetry at all. We also show that $T_1$ cannot be expressed by radicals (Proposition~\ref{prop:T1galois}) and that the three known solutions are isolated (Theorem~\ref{thm:localregular} and Section~\ref{sec:intervals}). The general case remains open. Possible approaches and open problems are discussed in Sections~\ref{tet:directions} and~\ref{sec:openproblems}.

\part{Counterexamples to the Acta converse conjectures}
\label{part:counterexamples}

We recall the construction in Section~\ref{sec:morley} and write the regularity of the Morley tetrahedron as a system of polynomial equations in Section~\ref{sec:gram}. Both counterexamples lie in a family of tetrahedra with two reflection symmetries, introduced in Section~\ref{sec:family}. The counterexamples are stated in Section~\ref{sec:counterexamples}, and the computations they rely on are given in Section~\ref{sec:intervals}.

\section{The Morley tetrahedron}\label{sec:morley}

Let $ABCD$ be a nondegenerate tetrahedron. For an edge $XY$, the two faces containing $XY$ form a dihedral angle. If $Z$ is one of the two vertices not on $XY$, we denote by
\[
(XY)_Z
\]
the trisecting plane through $XY$ which is closer to $Z$. Equivalently, if $\theta_{XY}$ is the interior dihedral angle along $XY$, then $(XY)_Z$ makes angle $\theta_{XY}/3$ with the face containing $XY$ and $Z$.

Following \cite{TranActa}, define the Morley tetrahedron
\[
\M(ABCD)=A'B'C'D'
\]
by
\begin{align}
A'&=(BC)_D\cap(CD)_B\cap(DB)_C, \notag\\
B'&=(AC)_D\cap(CD)_A\cap(DA)_C, \notag\\
C'&=(AB)_D\cap(BD)_A\cap(DA)_B, \notag\\
D'&=(AB)_C\cap(BC)_A\cap(CA)_B. \label{eq:morleydefinition}
\end{align}
Equivalently, $A'$ is the intersection of the three trisecting planes attached to the three edges of the face $BCD$ and lying closer to that face, and similarly for the other three vertices.

\begin{lemma}
\label{lem:equivariance} The Morley construction commutes with Euclidean isometries and with relabellings of the four vertices.
\end{lemma}

\begin{proof}
The construction uses only dihedral angles, their trisecting planes, the choice of the trisector closer to a specified face, and intersections of planes. All these operations commute with Euclidean isometries.

The description of a Morley vertex as the intersection of the three trisectors attached to the three edges of the opposite face is symmetric under relabelling.
\end{proof}

\section{The Gram matrix formulation}
\label{sec:gram}

Write the vertices as $A_1,A_2,A_3,A_4$, and let $F_i$ be the face opposite $A_i$. Denote its area by $S_i$, its inward unit normal by $n_i$, and the volume of the tetrahedron by $V$. If $\theta_{ij}$ is the interior dihedral angle between $F_i$ and $F_j$, set
\[
q_{ii}=1,\; q_{ij}=2\cos(\theta_{ij}/3)\;(i\ne j),
\; K=(q_{ij}),\; \ph(x)=\frac{3x-x^3}{2}.
\]
Thus $1<q_{ij}<2$ for $i\ne j$. Notice that $\theta_{ij}$ is the angle along the edge complementary to $A_iA_j$. The normal Gram matrix is
\begin{equation}
H=(n_i\cdot n_j)=\ph[K]=\frac{3K-K^{\circ3}}2,
\label{eq:normalgram}
\end{equation}
where the brackets and the superscript $\circ3$ indicate entrywise operations. The equilibrium of the face area vectors gives
\begin{equation}
H\succeq0,\; \rank H=3,\; HS=0,
\; S=(S_1,S_2,S_3,S_4)^T>0.
\label{eq:positivekernel}
\end{equation}

Let $\ell_j(X)$ be the inward signed distance to $F_j$, and let $M_i$ be the Morley vertex corresponding to $F_i$. The sine rule in a normal section gives
\begin{equation}
\ell_j(M_i)=t_iq_{ij},\;
t_i=\ell_i(M_i)=\frac{3V}{(KS)_i}>0.
\label{eq:distances}
\end{equation}
The distance ratio is $\sin(2\theta/3)/\sin(\theta/3)=2\cos(\theta/3)$. The second equality follows from the identity
\[
\sum_j S_j\ell_j(X)=3V.
\]
Conversely, the positive numbers in \eqref{eq:distances} are the distances of the interior point with barycentric coordinates $S_jt_iq_{ij}/(3V)$ to the faces. In particular, the three trisecting planes defining $M_i$ do meet in one point.

\begin{lemma}
\label{lem:negative}
The matrix $K$ is strictly negative definite on $\one^\perp$, where $\one=(1,1,1,1)^T$. Its inertia is $(1,3)$; in particular, it is invertible and $\det K<0$.
\end{lemma}

\begin{proof}
Define $f(h)=2\cos(\arccos(-h)/3)$. We claim that
\[
f(h)=\sqrt3-\sum_{m\ge1}c_mh^m,\;
c_m>0,\; c_1=\frac13,
\]
with absolute convergence on $[-1,1]$. Put $y=\sqrt3-f(h)$. The equation $\ph(f(h))=h$ becomes
\[
h=3y\left(1-\frac{y}{\sqrt3}\right)
       \left(1-\frac{y}{2\sqrt3}\right).
\]
Lagrange inversion gives
\[
c_m=\frac{1}{m3^m}[u^{m-1}]
\left(1-\frac{u}{\sqrt3}\right)^{-m}
\left(1-\frac{u}{2\sqrt3}\right)^{-m}>0.
\]
The branch is analytic for $|h|<1$, since the only possible finite branch values of the cubic are $h=\pm1$. Letting $h\uparrow1$ gives $\sum_m c_m=\sqrt3-1$, which proves the asserted endpoint convergence.

Each $H^{\circ m}$ is the Gram matrix of the tensors $n_i^{\otimes m}$ and is positive semidefinite. For $0\ne z\perp\one$, therefore,
\[
z^TKz=-\sum_{m\ge1}c_mz^TH^{\circ m}z
\le-\frac13z^THz<0.
\]
The last inequality follows from $\ker H=\R S$ and $S>0$. Finally, $\one^TK\one>0$, so the remaining eigenvalue is positive.
\end{proof}

\begin{proposition}[Regularity equations]
\label{prop:regularity}
If the Morley tetrahedron is regular with edge length $L$, set
\[
d_i=\frac{t_i}{\sqrt2L},\; d=(d_i)^T,
\; D=\diag(d_1,d_2,d_3,d_4).
\]
Then
\begin{equation}
H=K(4D^2-dd^T)K.
\label{eq:regularity}
\end{equation}
Conversely, suppose $K$ is symmetric with $q_{ii}=1$, $1<q_{ij}<2$ for $i\ne j$, and satisfies \eqref{eq:normalgram}, \eqref{eq:positivekernel} and \eqref{eq:regularity} for some $d>0$. The tetrahedron determined by $H$ has regular Morley tetrahedron.
\end{proposition}

\begin{proof}
Take the centroid of the regular Morley tetrahedron as origin, and write its vertices as $m_i$. Then
\[
\sum_i m_i=0,\;
m_i\cdot m_j=\frac{L^2}{2}\delta_{ij}-\frac{L^2}{8},
\; \sum_i m_im_i^T=\frac{L^2}{2}I_3.
\]
Using \eqref{eq:distances} gives
\[
n_j=\frac{2}{L^2}\sum_i t_iq_{ij}m_i.
\]
Taking scalar products proves \eqref{eq:regularity}.

For the converse, a positive semidefinite matrix $H$ with diagonal $1$, rank $3$, and a positive kernel vector determines four unit normals which positively span $\R^3$. The half spaces $n_i\cdot X\ge-1$ bound a nondegenerate tetrahedron with these inward normals; this is the simplex case of the standard facet-normal reconstruction, cf.~\cite{Schneider}. Its shape is uniquely determined by $H$: translations and a common dilation account for the four support parameters. The positivity of the kernel also implies that every three normals are independent.

By Lemma~\ref{lem:negative}, $K$ is invertible. Equation \eqref{eq:regularity} and the kernel condition imply
\[
KS=c(1/d_i)_i\;(c>0).
\]
Consequently the actual distances in \eqref{eq:distances} satisfy $t_i=\alpha d_i$ for a common $\alpha>0$. Put $T=\diag(t_i)$ and $B=K^{-1}HK^{-1}$. The Gram matrix of the barycentric gradients of the Morley tetrahedron is
\[
T^{-1}BT^{-1}=\alpha^{-2}(4I_4-\one\one^T).
\]
This is the Gram matrix of barycentric gradients of a regular tetrahedron, so all its edges are equal. The matrix of distance coordinates is invertible because $K$ and $T$ are invertible; thus the Morley tetrahedron is nondegenerate.
\end{proof}

In particular, regularity is equivalent to
\begin{equation}
\frac{B_{ij}}{\sqrt{B_{ii}B_{jj}}}=-\frac13\;(i\ne j).
\label{eq:normalizedgram}
\end{equation}
Indeed, \eqref{eq:normalizedgram} gives $B=4D^2-dd^T$ on taking $d_i=\sqrt{B_{ii}/3}$. Its polynomial form is $9B_{ij}^2=B_{ii}B_{jj}$ together with $B_{ij}<0$.

We shall repeatedly use the edge formula
\begin{equation}
|A_kA_l|=\frac{2S_iS_j\sin\theta_{ij}}{3V},
\; \{i,j,k,l\}=\{1,2,3,4\}.
\label{eq:edge}
\end{equation}
It follows by expressing the volume in terms of the two face altitudes to their common edge. Also,
\[
\sin^2\theta_{ij}=1-\ph(q_{ij})^2
=\frac{(q_{ij}^2-1)^2(4-q_{ij}^2)}4.
\]

\section{A symmetric two parameter family}\label{sec:family}

For $\rho>0$ and $0<\sigma<\pi/2$, let
\[
T(\rho,\sigma)=ABCD
\]
be given by
\[
A=\left(-\frac12,0,0\right), B=\left(\frac12,0,0\right),
\]
and
\[
C=(0,\rho\cos\sigma,\rho\sin\sigma), D=(0,\rho\cos\sigma,-\rho\sin\sigma).
\]
Its edge lengths are
\[
AB=1, AC=AD=BC=BD=\lambda=\sqrt{\frac14+\rho^2},
CD=\mu=2\rho\sin\sigma.
\]

The reflections
\[
(x,y,z)\longmapsto(-x,y,z)
\]
and
\[
(x,y,z)\longmapsto(x,y,-z)
\]
preserve the vertex set, so by Lemma~\ref{lem:equivariance} they are also symmetries of the Morley tetrahedron. Consequently
\[
|A'C'|=|B'C'|=|A'D'|=|B'D'|.
\]

Define
\[
f(\rho,\sigma)=|A'B'|^2-|C'D'|^2,
g(\rho,\sigma)=|A'B'|^2-|A'C'|^2.
\]
The Morley tetrahedron of $T(\rho,\sigma)$ is regular exactly when
\[
f(\rho,\sigma)=g(\rho,\sigma)=0,
\]
and, if it is not regular, it is isosceles exactly when
\[
f(\rho,\sigma)=0.
\]

\section{Counterexamples to the original conjectures}
\label{sec:counterexamples}

\begin{theorem}
\label{thm:T1} There exists a nondegenerate tetrahedron $T_1$ satisfying
\[
AB=1, AC=AD=BC=BD=\lambda_*,
CD=\mu_*,
\]
where
\[
\lambda_*\approx1.5793920543,
\]
and
\[
\mu_*\approx2.3518923337,
\]
whose Morley tetrahedron is regular. The tetrahedron $T_1$ is neither regular nor isosceles.
\end{theorem}

\begin{proof}
In Section~\ref{sec:intervals} we show that the regularity equations have exactly one solution in a small box around the values labelled $T_1$, and that this solution comes from a nondegenerate tetrahedron. The box and the equations are invariant under two disjoint transpositions, so the solution is invariant as well and the tetrahedron has the two corresponding reflections. After relabelling and scaling it belongs to the family $T(\rho,\sigma)$ of Section~\ref{sec:family}.

The formulas in Section~\ref{sec:intervals} give
\[
1.57<\lambda_*<1.59,\; 2.35<\mu_*<2.36.
\]
Hence $T_1$ is neither regular nor isosceles. The decimal values in the statement are approximations, and the parameters are
\[
\rho_*\approx1.4981586235,\; \sigma_*\approx0.9025793893.
\]
An exact algebraic specification of $T_1$ is given in Theorem~\ref{sym:thm:classification}.
\end{proof}

\subsection{A disphenoid counterexample}

For $c>0$, consider
\[
A=\left(-\frac12,0,0\right), B=\left(\frac12,0,0\right),
\]
and
\[
C=\left(0,c,\frac12\right), D=\left(0,c,-\frac12\right).
\]
Then
\[
AB=CD=1,
AC=AD=BC=BD= \sqrt{c^2+\frac12}.
\]
Thus every tetrahedron in this family is isosceles.

\begin{theorem}
\label{thm:T2} Let
\[
c_0^2=\frac{8\sqrt6-3}{90}.
\]
The isosceles tetrahedron $T_2$ of the family above with $c=c_0$ has a regular Morley tetrahedron. Its edges are
\[
AB=CD=1,\;
AC=AD=BC=BD=\lambda_\dagger=\sqrt{\frac{21+4\sqrt6}{45}}\approx0.8272841641,
\]
so $T_2$ is isosceles but not regular.
\end{theorem}

\begin{proof}
Put $r=\sqrt6$ and
\begin{equation}\label{t2:eq:exactpq}
q_0=2r-3,\;
p_0=r-\frac52+\frac12\sqrt{60r-135}.
\end{equation}
The number $p_0$ is the positive root of
\begin{equation}\label{t2:eq:pquadratic}
p^2+(5-2r)p+46-20r=0.
\end{equation}
The quadratic is negative at $1$ and positive at $2$, and its constant term is negative, so $1<p_0<2$. Also $1<q_0<2$ because $2<r<5/2$. Reducing powers of $p_0$ by \eqref{t2:eq:pquadratic} gives $p_0^3=3p_0+470-192r$, and a direct expansion gives $q_0^3=102r-243$. Hence
\begin{equation}\label{t2:eq:phi}
\ph(p_0)=96r-235,\; \ph(q_0)=117-48r.
\end{equation}

In the coordinates above with $c=c_0$, the inward unit normals of the faces opposite $A$ and $B$ have scalar product
\[
\frac{1/4-c_0^2}{1/4+c_0^2}=96\sqrt6-235=\ph(p_0),
\]
and the same holds for the faces opposite $C$ and $D$. For a face opposite $A$ or $B$ and a face opposite $C$ or $D$ the scalar product is
\[
-\frac{1/4}{1/4+c_0^2}=117-48\sqrt6=\ph(q_0).
\]
Since $\ph$ is strictly decreasing on $(1,2)$, formula \eqref{eq:normalgram} shows that the matrix $K$ of Section~\ref{sec:gram} is
\begin{equation}\label{t2:eq:K}
K=\begin{pmatrix}
1&p_0&q_0&q_0\\p_0&1&q_0&q_0\\q_0&q_0&1&p_0\\q_0&q_0&p_0&1
\end{pmatrix}.
\end{equation}
The row sums of $H=\ph[K]$ vanish, because
\begin{equation}\label{t2:eq:row}
1+\ph(p_0)+2\ph(q_0)=0.
\end{equation}
On $(1,1,-1,-1)^T$ the eigenvalues of $K$ and $H$ are $1+p_0-2q_0$ and $1+\ph(p_0)-2\ph(q_0)=-4\ph(q_0)$. On the plane spanned by $(1,-1,0,0)^T$ and $(0,0,1,-1)^T$ they are $1-p_0$ and $1-\ph(p_0)=(p_0-1)^2(p_0+2)/2$. Put $\tau_0^2=(p_0+2)/8$. Reduction by \eqref{t2:eq:pquadratic} gives
\begin{equation}\label{t2:eq:tau}
(p_0+2)(1+p_0-2q_0)^2=384r-936=-8\ph(q_0).
\end{equation}
Hence $H$ and $\tau_0^2K(4I_4-\one\one^T)K$ agree on all four eigenvectors, the constant vector being annihilated by both. This is \eqref{eq:regularity} with $d_1=\dots=d_4=\tau_0$. Proposition~\ref{prop:regularity} shows that the Morley tetrahedron is regular. Finally
\[
AC^2=c_0^2+\frac12=\frac{21+4\sqrt6}{45}\ne1.
\]
The interval argument in Section~\ref{sec:intervals} additionally proves that this solution is isolated among all tetrahedron shapes.
\end{proof}

\begin{theorem}
\label{thm:curve} There exists a real analytic one parameter family of nonisosceles tetrahedra whose Morley tetrahedra are isosceles and nonregular.
\end{theorem}

\begin{proof}
Use the equivalent parameters
\[
A=(-a,0,0),\ B=(a,0,0),\ C=(0,c,b),\ D=(0,c,-b),
\]
where $a=1/2$, $b=\rho\sin\sigma$ and $c=\rho\cos\sigma$. By Section~\ref{sec:intervals}, $\partial f/\partial b\ne0$ at $T_1$, so the implicit function theorem gives a real analytic curve $f=0$ through $T_1$.

Section~\ref{sec:intervals} also shows that no other tetrahedron close to $T_1$ has a regular Morley tetrahedron. Indeed, the dihedral parameters of the tetrahedron with parameters $(b,c)$ depend continuously on $(b,c)$, and so do the normalized distances $d_i$ when the Morley tetrahedron is regular; a regular example close to $T_1$ would therefore give a second zero in the uniqueness box. Hence $g\ne0$ at all points of the curve near $T_1$ other than $T_1$ itself, and the Morley tetrahedra at these points are isosceles but not regular. Finally, $AB\ne CD$ near $T_1$, so these tetrahedra are not isosceles.
\end{proof}

\begin{corollary}
The two converse conjectures proposed in \cite{TranActa} are false.
\end{corollary}

\begin{proof}
Theorem~\ref{thm:T1} disproves the regular case of Conjecture~\ref{conj:old1}, and Theorem~\ref{thm:curve} disproves its isosceles case. By Definition~4 and Remark~7 of \cite{TranActa}, the admissible choice $e=f=g=3$ makes every $p$ secting plane in the construction of Theorem~6 the corresponding trisecting plane. Conjecture~2 of \cite{TranActa} applies that construction to an arbitrary tetrahedron. Thus $\M_{3,3,3}(ABCD)=\M(ABCD)$ for every $ABCD$, and the same examples disprove both cases of Conjecture~\ref{conj:old2}.
\end{proof}

\begin{remark}
The direct edge-equality statement remains valid. Theorem~\ref{thm:T2} shows, however, that a nonregular isosceles tetrahedron can have a regular Morley tetrahedron. Thus the nonregularity in Definition~1 of \cite{TranActa} cannot be required of the output in Theorem~3 there. Theorem~\ref{thm:curve} supplies the counterexamples to the converse of the edge-equality statement.
\end{remark}

\section{Certified existence and local uniqueness}
\label{sec:intervals}

This section contains the computations used in Theorems~\ref{thm:T1}--\ref{thm:curve}. All decimal numbers below are exact rational numbers. The argument is the usual combination of the contraction mapping theorem with interval arithmetic; see \cite{Krawczyk,Neumaier,Rump}.

\subsection{The equations and the three centers}

Use the coordinate order
\[
x=(q_{12},q_{13},q_{14},q_{23},q_{24},q_{34},d_1,d_2,d_3,d_4).
\]
Set $u=Kd$ and, for $i\le j$, define
\begin{equation}
F_{ij}(x)=4\sum_{k=1}^4d_k^2q_{ik}q_{jk}-u_i u_j-\ph(q_{ij}).
\label{eq:tenequations}
\end{equation}
Thus $F(x)=0$ is precisely \eqref{eq:regularity}. Order its components as
\[
(F_{11},F_{12},F_{13},F_{14},F_{22},F_{23},F_{24},F_{33},F_{34},F_{44}).
\]
The three rational centers have the common form
\begin{equation}
x_0=(u_0,u_0,v_0,w_0,u_0,u_0,a_0,b_0,b_0,a_0),
\label{eq:centerform}
\end{equation}
with the entries given in the following table.

\begin{center}
\small
\begin{tabular}{lrrrrr}
\toprule
 & $u_0$ & $v_0$ & $w_0$ & $a_0$ & $b_0$\\
\midrule
Regular & 1.833986596704 & 1.833986596704 & 1.833986596704 & 0.692277635482 & 0.692277635482\\
$T_1$ & 1.844867969052 & 1.648726879892 & 1.892477377433 & 0.675344993308 & 0.697538294418\\
$T_2$ & 1.898979485567 & 1.679329660736 & 1.679329660736 & 0.678171222916 & 0.678171222916\\
\bottomrule
\end{tabular}
\end{center}

Around each center take the closed box
\[
\mathcal I=\{x:\norm{x-x_0}_\infty\le r\},
\; r=10^{-5}.
\]
The matrix $J_0=DF(x_0)$ is nonsingular at all three centers. Define
\[
C=J_0^{-1}.
\]
Its entries are rational numbers determined by \eqref{eq:tenequations} and the table.

Here are the derivative formulas used in the calculation. For $a<b$, let $E^{ab}$ have entries $1$ in positions $(a,b),(b,a)$ and zero elsewhere. Then
\begin{align}
\frac{\partial F_{ij}}{\partial q_{ab}}
={}&4\sum_k d_k^2\bigl(E^{ab}_{ik}q_{jk}+q_{ik}E^{ab}_{jk}\bigr)
-(E^{ab}d)_i u_j-u_i(E^{ab}d)_j
-\ph'(q_{ij})E^{ab}_{ij},\notag\\
\frac{\partial F_{ij}}{\partial d_a}
={}&8d_aq_{ia}q_{ja}-q_{ia}u_j-u_iq_{ja}.
\label{eq:tenjacobian}
\end{align}
These formulas give $J_0$ by rational substitution. On substituting the intervals $[x_{0,k}-r,x_{0,k}+r]$ into \eqref{eq:tenjacobian}, ordinary interval addition and multiplication give a matrix of intervals $J(\mathcal I)$ containing every $DF(x)$ for $x\in\mathcal I$. For all three rows of the table, this calculation gives
\begin{equation}
\norm{I-CJ(\mathcal I)}_\infty<\frac7{100},
\; \norm{CF(x_0)}_\infty<10^{-11}.
\label{eq:contractionbounds}
\end{equation}
For a matrix of intervals, the norm is the largest row sum of the largest absolute endpoints of the entries. Thus \eqref{eq:contractionbounds} is a finite list of comparisons between rational numbers determined by the table and the formulas above.

The map $x\mapsto x-CF(x)$ is a contraction on $\mathcal I$, and its displacement from the center is at most
\[
10^{-11}+\frac7{100}r<r.
\]
It therefore has a unique fixed point $x_*$ in $\mathcal I$, and since $C$ is invertible, $x_*$ is the unique zero of $F$ in $\mathcal I$. Moreover
\begin{equation}
\norm{x_*-x_0}_\infty\le\frac{\norm{CF(x_0)}_\infty}{1-7/100}<1.1\cdot10^{-11}.
\label{eq:sharpbox}
\end{equation}

The same interval operations, using the determinant formula and Cramer's rule, give throughout each box
\begin{equation}
1<q_{ij}<2,\; d_i>0,\; \det K<0,
\; K^{-1}(1/d_i)_i>0.
\label{eq:intervalgeometry}
\end{equation}
At the zero of $F$, the matrix $H=K(4D^2-dd^T)K$ is positive semidefinite of rank $3$. The last vector in \eqref{eq:intervalgeometry} belongs to its kernel and has positive coordinates. Proposition~\ref{prop:regularity} therefore gives a nondegenerate tetrahedron whose Morley tetrahedron is regular.

Every permutation of the vertices preserves the regular center. The centers of $T_1$ and $T_2$ are preserved by $(1\,4)$ and $(2\,3)$, and that of $T_2$ also by $(1\,2)(3\,4)$. These permutations preserve the boxes and the equations, so the unique solution in each box has the same symmetries. In particular, up to similarity and relabelling, each of the three tetrahedra is the only example in its neighbourhood.

\subsection{The edge lengths and the analytic curve}

Relabel the two reflections as $(1\,2),(3\,4)$, and put $p=q_{12}$ and $s=q_{34}$. Denote the common value of the four cross entries by $q$, and set $h_p=\ph(p)$ and $h_s=\ph(s)$. The original tetrahedron has coordinates
\[
A=(-a,0,0),\; B=(a,0,0),\;
C=(0,c,b),\; D=(0,c,-b),
\]
where $a=1/2$ and
\begin{equation}
c^2=a^2\frac{1-h_p}{1+h_p},\;
b^2=c^2\frac{1+h_s}{1-h_s}.
\label{eq:reconstruction}
\end{equation}
For $T_1$, this means $p=q_{23}$ and $s=q_{14}$ in \eqref{eq:centerform}; for $T_2$, take $p=q_{14}=q_{23}$. The edge $CD$ has length $2b$, and each cross edge has length $\sqrt{a^2+b^2+c^2}$. By \eqref{eq:sharpbox}, the true values of $p$ and $s$ lie within $1.1\cdot10^{-11}$ of the entries of the table. Substituting these intervals gives
\[
1.57<\lambda_*<1.59,\; 2.35<\mu_*<2.36
\]
for $T_1$, and
\[
0.82<\lambda_\dagger<0.84
\]
for $T_2$. In the latter case $p=s$, and \eqref{eq:reconstruction} gives $b=a$. These are the bounds used in Theorem~\ref{thm:T1}; for $T_2$ they agree with the exact value in Theorem~\ref{thm:T2}.

We finally check the derivative used in Theorem~\ref{thm:curve}. Put
\[
U=(a^2+c^2)^{1/2},\; W=(b^2+c^2)^{1/2}.
\]
The two distinct face areas are $bU$ and $aW$. Formula \eqref{eq:distances} and the two reflections give
\[
\begin{aligned}
t_A&=\frac{2abc}{(1+p)bU+2qaW},&
t_C&=\frac{2abc}{2qbU+(1+s)aW},\\
|A'B'|^2&=\frac{(p-1)^2t_A^2U^2}{c^2},&
|C'D'|^2&=\frac{(s-1)^2t_C^2W^2}{c^2}.
\end{aligned}
\]
With $a,c$ fixed, the normal Gram entries satisfy
\[
h_p=\frac{a^2-c^2}{a^2+c^2},\;
h_s=\frac{b^2-c^2}{b^2+c^2},\;
h_q=-\frac{ab}{UW}.
\]
Hence
\[
p_b=0,\;
s_b=\frac{4bc^2}{(b^2+c^2)^2\ph'(s)},\;
q_b=-\frac{ac^2}{UW^3\ph'(q)}.
\]
We differentiate the formulas for $|A'B'|^2$ and $|C'D'|^2$ with respect to $b$ and substitute the intervals given by \eqref{eq:sharpbox} for $T_1$. This gives
\[
-0.001840<\frac{\partial f}{\partial b}<-0.001839.
\]
Only rational operations and square roots occur, and a square root of a positive rational interval is enclosed by rational endpoints whose squares bound the interval. In particular $\partial f/\partial b\ne0$, as needed in the proof of Theorem~\ref{thm:curve}.

\part{New conjectures}
\label{part:conjectures}

The counterexamples are not arbitrary. Both $T_1$ and $T_2$ have four equal cross edges, or equivalently two perpendicular reflection planes interchanging disjoint pairs of vertices. This suggests a weaker converse, in which regularity of the Morley tetrahedron forces some symmetry of the original tetrahedron but not regularity. In this part we state this and related conjectures. The partial results are proved in Part~\ref{part:results}.

\section{What the counterexamples suggest}
\label{sec:suggest}

In both exceptional examples there is a pair of opposite edges, $AB$ and $CD$, such that the other four edges have a common length. The plane through $C$, $D$ and the midpoint of $AB$ is then the perpendicular bisector plane of $AB$, and the reflection in it interchanges $A$ and $B$ and fixes $C$ and $D$. In the same way the perpendicular bisector plane of $CD$ contains $A$ and $B$. The two planes are perpendicular, because $AB\perp CD$. Conversely, two such reflections force the four cross edges to be equal. In $T_2$ the edges $AB$ and $CD$ are equal as well; in $T_1$ they are not.

Section~\ref{sec:intervals} proves existence and uniqueness near the three known examples, but says nothing about the rest of the shape space.

\section{The conjectures}
\label{sec:conjectures}

\begin{conjecture}[Four equal edges conjecture]
\label{conj:four} Let $ABCD$ be a nondegenerate tetrahedron. If its Morley tetrahedron is regular, then, after relabelling the vertices,
\[
AC=AD=BC=BD.
\]
In other words, there is a pair of opposite edges such that the remaining four edges, which form a skew quadrilateral, all have the same length. Equivalently, the tetrahedron has two reflection symmetries interchanging the endpoints of these two opposite edges.
\end{conjecture}

These examples suggest a complete classification.

\begin{conjecture}[Classification]\label{conj:classification}
Up to similarity and relabelling, the only tetrahedra with regular Morley tetrahedron are the regular tetrahedron and the two exceptional similarity classes represented by $T_1$ and $T_2$.
\end{conjecture}

A tetrahedron with four equal cross edges has two reflection symmetries. The next conjecture asks only for one symmetry of any kind.

\begin{conjecture}[Symmetry conjecture]\label{conj:symmetry}
If the Morley tetrahedron of $ABCD$ is regular, then some isometry other than the identity maps the vertex set $\{A,B,C,D\}$ onto itself.
\end{conjecture}

\begin{conjecture}[One pair of equal opposite edges]
\label{conj:onepair}
Let $ABCD$ be a nondegenerate tetrahedron satisfying $AB=CD$. If its Morley tetrahedron is regular, then, after a possible relabelling,
\[
AC=AD=BC=BD.
\]
\end{conjecture}

Conjecture~\ref{conj:onepair} is a special case of Conjecture~\ref{conj:four}. Note that its conclusion does not say that the remaining two edges are equal: after relabelling, $T_1$ satisfies $AB=CD$ and has four equal cross edges, but it is not isosceles.

In the notation of Section~\ref{sec:gram}, Conjecture~\ref{conj:four} can be stated purely in terms of matrices.

\begin{conjecture}[Positivity lemma]
\label{conj:positivity}
Let $K$ be symmetric with $q_{ii}=1$ and $1<q_{ij}<2$ for $i\ne j$. Suppose that $H=\ph[K]$ satisfies \eqref{eq:positivekernel}, and that $B=K^{-1}HK^{-1}$ satisfies \eqref{eq:normalizedgram}. Then, after a permutation of the indices,
\[
q_{13}=q_{14}=q_{23}=q_{24}.
\]
\end{conjecture}

Conjecture~\ref{conj:positivity} is equivalent to Conjecture~\ref{conj:four}. Indeed, the displayed equalities make $H$ invariant under $(1\,2)$ and $(3\,4)$, so its positive kernel vector satisfies $S_1=S_2$ and $S_3=S_4$, and formula \eqref{eq:edge} gives four equal cross edges. Conversely, four equal cross edges give two reflections, and the displayed equalities follow by equivariance. Proposition~\ref{prop:regularity} provides the passage from matrices to tetrahedra.

Finally, we state the analogue for the sector construction of \cite{TranActa}.

\begin{conjecture}[$p$ sector four equal edges conjecture]
\label{conj:psector} Let $\M_{e,f,g}(ABCD)$ denote the generalized $p$ sector tetrahedron of \cite[Theorem 6]{TranActa}, for admissible parameters $e,f,g>1$. If $\M_{e,f,g}(ABCD)$ is regular, then $ABCD$ has, after relabelling, four equal edges.
\end{conjecture}

\subsection*{Relations between the conjectures}

Clearly Conjecture~\ref{conj:four} implies Conjecture~\ref{conj:onepair}. The results of Part~\ref{part:results} show more:
\begin{itemize}
\item Conjectures~\ref{conj:four} and~\ref{conj:classification} are equivalent (Corollary~\ref{sym:cor:equivalence});
\item Conjectures~\ref{conj:four} and~\ref{conj:symmetry} are equivalent (Corollary~\ref{cor:symmetric});
\item Conjectures~\ref{conj:four} and~\ref{conj:positivity} are equivalent (see above);
\item under $AB=CD$, Conjecture~\ref{conj:onepair} is equivalent to the equality of the dihedral angles along $AB$ and $CD$ (Proposition~\ref{prop:cotangent}).
\end{itemize}
In particular, to prove Conjecture~\ref{conj:four} it is enough to find one reflection symmetry; Theorem~\ref{thm:reflection} then gives the second one.

\part{Results, possible approaches, and open problems}
\label{part:results}

In this part we prove the conjectures of Part~\ref{part:conjectures} under additional assumptions. Section~\ref{sec:bounds} contains a priori bounds and a criterion for reflection symmetry. Isosceles tetrahedra are treated in Section~\ref{sec:isosceles}. Section~\ref{sec:symmetry} deals with two pairs of equal opposite edges, with a single reflection, with four equal cross edges, and finally with an arbitrary symmetry. The case $AB=CD$ is reformulated in Section~\ref{sec:onepair}, and Section~\ref{sec:nearregular} proves local uniqueness at the regular tetrahedron. We summarize the situation in Section~\ref{sec:summary} and discuss possible approaches and open problems in Sections~\ref{tet:directions} and~\ref{sec:openproblems}.

\section{Bounds and a reflection criterion}
\label{sec:bounds}

We begin with bounds on the distances $t_i$ of \eqref{eq:distances}.

\begin{proposition}
\label{prop:distancebounds}
If $r_{\mathrm{in}}$ is the inradius of the original tetrahedron, then
\[
\frac{r_{\mathrm{in}}}{\sqrt3}<t_i<\frac{2r_{\mathrm{in}}}{3},
\; \frac{\max_i t_i}{\min_i t_i}<\frac2{\sqrt3}.
\]
These inequalities do not require the Morley tetrahedron to be regular.
\end{proposition}

\begin{proof}
Normalize the face areas to sum to $1$ and put $y_i=(KS)_i$. Since $HS=0$, we have $\sum_jS_jq_{ij}^3=3y_i$. Strict convexity gives $y_i^3<3y_i$, hence $y_i<\sqrt3$. On the other hand,
\[
x^3-7x+6=(x-1)(x-2)(x+3)<0\;(1<x<2).
\]
Applying this to the off-diagonal terms gives $3y_i<7y_i-6$, so $y_i>3/2$. Now $t_i=r_{\mathrm{in}}/y_i$.
\end{proof}

\begin{proposition}[A reflection criterion]
\label{prop:defect}
If the Morley tetrahedron is regular with edge length $L$, then
\begin{equation}
\varepsilon_{ij}:=q_{ij}+2-\frac{(t_i+t_j)^2}{L^2}\ge0.
\label{eq:defect}
\end{equation}
Equality holds if and only if the original tetrahedron has a reflection interchanging $A_i,A_j$ and fixing the other two vertices. Moreover,
\begin{equation}
\frac{(t_i+t_j)^2}{L^2(q_{ij}+2)}
+\frac{(t_i-t_j)^2}{L^2(2-q_{ij})}\le1.
\label{eq:ellipse}
\end{equation}
In particular, $d_i\le1/\sqrt2$.
\end{proposition}

\begin{proof}
Set $v=M_j-M_i$ and $q=q_{ij}$. Formula \eqref{eq:distances} gives
\[
(n_i-n_j)\cdot v=(q-1)(t_i+t_j),
\; (n_i+n_j)\cdot v=(q+1)(t_j-t_i).
\]
The two normal sums are orthogonal and have squared lengths
\[
\|n_i-n_j\|^2=(q-1)^2(q+2),
\; \|n_i+n_j\|^2=(2-q)(q+1)^2.
\]
Cauchy's inequality proves \eqref{eq:defect}; summing the squared orthogonal projections of $v$ proves \eqref{eq:ellipse}. Minimizing its left side with respect to $t_j$ gives $t_i^2/L^2\le1$, or $d_i\le1/\sqrt2$.

If $\varepsilon_{ij}=0$, then $v$ is parallel to $n_i-n_j$, and the second scalar product implies $t_i=t_j$. The plane $\ell_i=\ell_j$ is the perpendicular bisector of $M_iM_j$. The other two Morley vertices lie in that plane, whence $q_{ik}=q_{jk}$ for $k\ne i,j$. Also $n_k\cdot v=0$ for these $k$. Reflection in the plane therefore interchanges $F_i,F_j$ and preserves the other two faces. It induces exactly the stated permutation of the original vertices. Conversely, such a reflection is preserved by the Morley construction, and $v$ is parallel to $n_i-n_j$, giving equality in Cauchy's inequality.
\end{proof}

When equality holds, the distances $t_i$ and $t_j$ can be computed.

\begin{corollary}\label{cor:reflectiondistance}
Suppose that the Morley tetrahedron is regular with edge length $L$. If the original tetrahedron has a reflection interchanging $A_i,A_j$ and fixing the other two vertices, then
\[
t_i=t_j=L\cos\frac{\theta_{ij}}6,
\;\text{equivalently}\;
d_i^2=d_j^2=\frac{q_{ij}+2}8 .
\]
\end{corollary}

\begin{proof}
By Proposition~\ref{prop:defect}, $\varepsilon_{ij}=0$ and $t_i=t_j$. Hence $4t_i^2=L^2(q_{ij}+2)$. Since $q_{ij}+2=2\cos(\theta_{ij}/3)+2=4\cos^2(\theta_{ij}/6)$, the claim follows.
\end{proof}

The relation $8a^2=p+2$ in the proofs of Theorems~\ref{t2:thm:classification}, \ref{thm:reflection} and~\ref{sym:thm:classification} is this identity.

\section{The isosceles case and equal face distances}
\label{sec:isosceles}

\begin{theorem}
\label{thm:disphenoid}
Let $ABCD$ be an isosceles tetrahedron. If its Morley tetrahedron is regular, then, after relabelling,
\[
AC=AD=BC=BD.
\]
Thus at least two of the three common lengths of opposite edge pairs coincide. In this case the remaining two edges also have equal lengths, because the original tetrahedron is isosceles.
\end{theorem}

\begin{proof}
All four faces of an isosceles tetrahedron are congruent. The three rotations through $\pi$ that permute the vertices give
\[
K=\begin{pmatrix}
1&a&b&c\\ a&1&c&b\\ b&c&1&a\\ c&b&a&1
\end{pmatrix},\; S_1=S_2=S_3=S_4,
\]
where $a,b,c>1$. Put $\sigma=1+a+b+c$. By \eqref{eq:distances}, all $t_i$ are equal, so all $d_i$ equal some $\tau>0$.

On the vectors $(1,1,-1,-1)^T$, $(1,-1,1,-1)^T$, and $(1,-1,-1,1)^T$, the eigenvalues of $K$ are respectively $2a+2-\sigma$, $2b+2-\sigma$, and $2c+2-\sigma$. Since $H\one=0$, the corresponding eigenvalues of $H$ are $2+3a-a^3$, $2+3b-b^3$, and $2+3c-c^3$. Equation \eqref{eq:regularity} therefore gives
\[
2+3u-u^3=4\tau^2(2u+2-\sigma)^2
\;(u=a,b,c).
\]
If $a,b,c$ were distinct, they would be all three roots of the monic cubic
\[
X^3-3X-2+4\tau^2(2X+2-\sigma)^2.
\]
Vi\`ete's formula would give $a+b+c=-16\tau^2<0$, contrary to $a,b,c>1$. Hence two of $a,b,c$ coincide. Together with the equality of the face areas, formula \eqref{eq:edge} now gives four equal cross edges.
\end{proof}

In fact it suffices to know that the four distances $t_i$ are equal.

\begin{theorem}
\label{thm:equaldistances}
Suppose the Morley tetrahedron is regular and $t_1=t_2=t_3=t_4$. Then the original tetrahedron is isosceles and has four equal cross edges after relabelling.
\end{theorem}

\begin{proof}
Write
\[
\begin{aligned}
q_{12}&=a+x,& q_{34}&=a-x,\\
q_{13}&=b+y,& q_{24}&=b-y,\\
q_{14}&=c+z,& q_{23}&=c-z.
\end{aligned}
\]
Here $a,b,c>1$. In the orthonormal Hadamard basis, with the constant vector first, $K$ becomes
\[
\widehat K=\begin{pmatrix}
\lambda_0&x&y&z\\
x&\lambda_1&-z&-y\\
y&-z&\lambda_2&-x\\
z&-y&-x&\lambda_3
\end{pmatrix},
\]
where $\lambda_0=1+a+b+c$ and
\[
\lambda_1=1+a-b-c,\;
\lambda_2=1-a+b-c,\;
\lambda_3=1-a-b+c.
\]
By Lemma~\ref{lem:negative}, $\mu_i=-\lambda_i>0$ and
\begin{equation}
A=\begin{pmatrix}\mu_1&z&y\\z&\mu_2&x\\y&x&\mu_3\end{pmatrix}\succ0.
\label{eq:equalA}
\end{equation}

Let $d_i=\tau>0$. In this basis, \eqref{eq:regularity} reads
\[
\widehat H=4\tau^2\widehat K\diag(0,1,1,1)\widehat K.
\]
The four diagonal entries of $H$ are equal. Equivalently,
\[
\widehat H_{01}+\widehat H_{23}
=\widehat H_{02}+\widehat H_{13}
=\widehat H_{03}+\widehat H_{12}=0.
\]
Multiplication gives
\begin{equation}
\begin{aligned}
(\mu_2+\mu_3-\mu_1)x&=yz,\\
(\mu_1+\mu_3-\mu_2)y&=xz,\\
(\mu_1+\mu_2-\mu_3)z&=xy.
\end{aligned}
\label{eq:equalxyz}
\end{equation}

Suppose first that $xyz\ne0$, and denote the three coefficients on the left by $\eta_1,\eta_2,\eta_3$. They have the same sign, and their sum is positive, so all are positive. Also
\[
x^2=\eta_2\eta_3,\; y^2=\eta_1\eta_3,\;
z^2=\eta_1\eta_2,\; xyz=\eta_1\eta_2\eta_3.
\]
Since $\mu_1=(\eta_2+\eta_3)/2$ and cyclically,
\[
\det A=-\frac38\left(\sum_{i\ne j}\eta_i^2\eta_j
-6\eta_1\eta_2\eta_3\right)\le0
\]
by the inequality between the arithmetic and geometric means. This contradicts \eqref{eq:equalA}. Thus one of $x,y,z$ vanishes, and \eqref{eq:equalxyz} implies that at least two vanish. Relabel so that $x=y=0$.

Assume $z\ne0$. The last equation in \eqref{eq:equalxyz} gives
\begin{equation}
3c=a+b+1.
\label{eq:equalc}
\end{equation}
If $a\ne b$, comparison of $\widehat H_{11}-\widehat H_{22}$ and $\widehat H_{03}$, using both \eqref{eq:normalgram} and \eqref{eq:regularity}, gives
\[
a^2+ab+b^2=3c^2+z^2.
\]
Put $m=(a+b)/2>1$ and $e=(a-b)/2$. The preceding equality and \eqref{eq:equalc} yield
\[
z^2-\lambda_1\lambda_2
=\frac{(m-1)(11m+7)}9+5e^2>0,
\]
contradicting the positive principal minor $\lambda_1\lambda_2-z^2$ in \eqref{eq:equalA}.

It remains to consider $a=b$. Set $r=c-1>0$ and $w=z^2/r^2>0$. Then
\[
a=b=1+\frac32r,\;
\lambda_1=\lambda_2=-r,\; \lambda_3=-2r.
\]
The trace equation gives $\tau^2=1/[3r^2(2+w)]$. Comparison of $\widehat H_{00}-\widehat H_{33}$ and of $\widehat H_{03}$ gives respectively
\[
r^2(r+2)=\frac{32(5+w)}{81(2+w)},
\;
3r^2(2+w)\bigl(6+r(3+w)\bigr)=16.
\]
Eliminating $r^2$ gives
\[
r(2w^2+16w+3)+12w+6=0,
\]
which is impossible. Therefore $z=0$ as well.

We have proved equality of the three pairs of opposite entries of $K$. Its row sums are constant, and \eqref{eq:regularity} gives $H\one=0$. Since the kernel of $H$ has dimension one, the four face areas are equal. Formula \eqref{eq:edge} proves that the original tetrahedron is isosceles, and Theorem~\ref{thm:disphenoid} completes the proof.
\end{proof}

\subsection{Exact disphenoid classification}
\begin{theorem}\label{t2:thm:classification}
Up to similarity and relabelling, exactly two isosceles tetrahedra have
a regular Morley tetrahedron: the regular tetrahedron and the tetrahedron
$T_2$ with
\[
 AB=CD=1,\;
 AC=AD=BC=BD=\sqrt{\frac{21+4\sqrt6}{45}}.
\]
Here \emph{isosceles} means that all three pairs of opposite edges
have equal lengths.
\end{theorem}
\begin{proof}
By Theorem~\ref{thm:disphenoid}, after relabelling we may write
\begin{equation}\label{t2c:eq:K}
 K=\begin{pmatrix}
 1&p&q&q\\p&1&q&q\\q&q&1&p\\q&q&p&1
 \end{pmatrix},\; 1<p,q<2.
\end{equation}

The kernel condition is
\begin{equation}\label{t2c:eq:row}
 1+\ph(p)+2\ph(q)=0.
\end{equation}
On $(1,1,-1,-1)^T$, the eigenvalues of $K,H$ are
$1+p-2q,-4\ph(q)$. On the two-dimensional subspace spanned by
$(1,-1,0,0)^T,(0,0,1,-1)^T$, they are $1-p,1-\ph(p)$.
Since
\[
 1-\ph(p)=\frac{(p-1)^2(p+2)}2,
\]
regularity forces
\begin{equation}\label{t2c:eq:tau}
 \tau^2=\frac{p+2}{8},\;
 (p+2)(1+p-2q)^2=-8\ph(q).
\end{equation}

\medskip\noindent\emph{Elimination.}

Equations \eqref{t2c:eq:row} through \eqref{t2c:eq:tau} give
\begin{align*}
 A(p,q)&=p^3+2q^3-3p-6q-2=0,\\
 B(p,q)&=(p+2)(p+1-2q)^2+12q-4q^3=0.
\end{align*}
Their resultant in $p$ has the exact factorization
\begin{equation}\label{t2:eq:resultant}
 \Res_p(A,B)
 =-8q^2(q^2+6q-15)^2(3q^3-9q-2).
\end{equation}
It can be checked by computing the determinant of the Sylvester
matrix of the two cubics (script \texttt{check\_T2.py}).

Since $1<q<2$, there are only two possibilities.
If $3q^3-9q-2=0$, then $\ph(q)=-1/3$, and
\eqref{t2c:eq:row} gives $\ph(p)=-1/3$. The function $\ph$ is strictly
decreasing on $(1,2)$, so $p=q$; the original tetrahedron is regular.

Otherwise $q^2+6q-15=0$, and its only root in $(1,2)$ is $q_0=2\sqrt6-3$.
By \eqref{t2c:eq:row} and the strict monotonicity of $\ph$, the value
of $p$ is then uniquely determined; by \eqref{t2:eq:phi} it is the
number $p_0$ of \eqref{t2:eq:exactpq}. The proof of
Theorem~\ref{thm:T2} shows that this pair is realized by the
tetrahedron $T_2$. This completes the proof.
\end{proof}

\section{Symmetry results}\label{sec:symmetry}

\begin{theorem}[Two pairs of equal opposite edges]
\label{thm:twopairs}
Suppose the Morley tetrahedron of a nondegenerate tetrahedron is regular. If the original tetrahedron has two pairs of equal opposite edges, then it has four equal cross edges after relabelling.
\end{theorem}

\begin{proof}
Relabel so that $A_1A_3=A_2A_4$ and $A_1A_4=A_2A_3$. The permutation $(1\,2)(3\,4)$ preserves all six edge lengths, hence is induced by an isometry. Equivariance gives
\begin{equation}
K=\begin{pmatrix}1&p&q&r\\p&1&r&q\\q&r&1&s\\r&q&s&1\end{pmatrix},
\; d=(a,a,b,b)^T,\; S=(\alpha,\alpha,\beta,\beta)^T,
\label{eq:twopairsK}
\end{equation}
with $1<p,q,r,s<2$ and $a,b,\alpha,\beta>0$. On the plus and minus eigenspaces of $(1\,2)(3\,4)$, the blocks of $K$ are
\[
K_+=\begin{pmatrix}1+p&q+r\\q+r&1+s\end{pmatrix},
\; K_-=\begin{pmatrix}1-p&q-r\\q-r&1-s\end{pmatrix}.
\]
Equation \eqref{eq:regularity} becomes
\begin{equation}
H_+=2K_+\binom{a}{-b}(a,-b)K_+,
\; H_-=4K_-\diag(a^2,b^2)K_-.
\label{eq:twopairsblocks}
\end{equation}

We show that $p=s$ or $q=r$. Otherwise, interchange the two vertex pairs so that $p>s$, and put
\[
P=p-1,\; R=s-1,\; h=q-r,\; k=q+r,
\; u=P+R,\; v=PR,\; w=\frac{h^2}{PR}.
\]
Then $0<R<P<1$, $h\ne0$, and $h^2<PR$ by Lemma~\ref{lem:negative}. The diagonal entries of the minus equation in \eqref{eq:twopairsblocks} give
\[
P^2(P+3)=8a^2P^2+8b^2h^2,
\; R^2(R+3)=8b^2R^2+8a^2h^2.
\]
Solving,
\begin{align}
8a^2&=\frac{R^2\{P^2(P+3)-h^2(R+3)\}}{P^2R^2-h^4},\notag\\
8b^2&=\frac{P^2\{R^2(R+3)-h^2(P+3)\}}{P^2R^2-h^4}.
\label{eq:twopairsab}
\end{align}
In particular,
\begin{equation}
0<w<\frac{R(R+3)}{P(P+3)}<1.
\label{eq:wbound}
\end{equation}
The off-diagonal minus equation, divided by $h$, gives
\begin{equation}
k^2=4+\frac{4(u^2+3u-2v)}{3(1+w)}-\frac{vw}{3}.
\label{eq:ksquared}
\end{equation}
The two diagonal plus equations are
\[
[a(P+2)-bk]^2=\frac{(1-P)(P+2)^2}{4},
\;
[ak-b(R+2)]^2=\frac{(1-R)(R+2)^2}{4}.
\]
Divide by $P+2$ and $R+2$ respectively, subtract to cancel $ab$, and use \eqref{eq:twopairsab}, \eqref{eq:wbound}, and \eqref{eq:ksquared}. After dividing by $P-R\ne0$, the resulting equation is
\begin{equation}
\begin{aligned}
0={}&[2u(2u+3-3w)-v(w^2+4w+11)]\\
&\cdot[v(u+5+w(u+1))+2wu(u+3)]\\
&-6(u+1)(v+2u+4)v(1-w)(1+w)^2.
\end{aligned}
\label{eq:Ezero}
\end{equation}

The positive kernel vector $(\alpha,\beta)^T$ of $H_+$, whose diagonal entries are positive, forces $\ph(q)+\ph(r)<0$. Since
\[
\ph(q)+\ph(r)=\frac{k(12-k^2-3h^2)}8,
\]
equation \eqref{eq:ksquared} makes this equivalent to
\begin{equation}
\mathcal B(u,v,w):=u^2+3u-6-6w+2v(w^2+w-1)>0.
\label{eq:Bpositive}
\end{equation}
The polynomial inequality in Appendix~\ref{app:bernstein} shows that \eqref{eq:wbound}, \eqref{eq:Ezero}, and \eqref{eq:Bpositive} are incompatible. Thus $p=s$ or $q=r$.

If $q=r$, formula \eqref{eq:edge} and the paired areas in \eqref{eq:twopairsK} immediately give the equality of the four cross edges. If $p=s$, the row sums of $H$ are equal. Their common value is zero, because $HS=0$ and $\one^TS>0$. Hence $S$ is constant, and \eqref{eq:edge} makes the tetrahedron isosceles. Theorem~\ref{thm:disphenoid} applies.
\end{proof}

\begin{theorem}[A single reflection]
\label{thm:reflection}
Suppose the Morley tetrahedron is regular. If the original tetrahedron has a reflection interchanging $A_1,A_2$ and fixing $A_3,A_4$, then it also has a reflection interchanging $A_3,A_4$ and fixing $A_1,A_2$. Consequently,
\[
A_1A_3=A_1A_4=A_2A_3=A_2A_4.
\]
\end{theorem}

\begin{proof}
The first reflection gives
\[
K=\begin{pmatrix}1&p&q&r\\p&1&q&r\\q&q&1&s\\r&r&s&1\end{pmatrix},
\; d=(a,a,b,c)^T.
\]
Appendix~\ref{app:reflection} excludes $b\ne c$ by exact polynomial elimination and the positivity of the face areas. If $b=c$, the identity
\[
(n_3+n_4)\cdot(n_3-n_4)=0
\]
in the coordinates used there gives $(q+r-(1+s)b/a)(q-r)=0$. The first factor cannot vanish: it would make $n_3+n_4=0$, contrary to $1<s<2$. Thus $q=r$. The matrix $H$ and its positive kernel are then also invariant under $(3\,4)$. The corresponding vertex permutation is an isometry by \eqref{eq:edge}, and is a reflection because it fixes two vertices and interchanges the other two.
\end{proof}

In terms of Proposition~\ref{prop:defect}, Theorem~\ref{thm:reflection} says that $\varepsilon_{12}=0$ implies $\varepsilon_{34}=0$. To prove Conjecture~\ref{conj:four} it would therefore be enough to show that one of the six defects $\varepsilon_{ij}$ vanishes. Note also that neither Theorem~\ref{thm:twopairs} nor Theorem~\ref{thm:reflection} covers the case of only one pair of equal opposite edges.

\subsection{Classification with four equal cross edges}\label{sym:section}

\begin{theorem}
\label{sym:thm:classification}
Suppose that $AC=AD=BC=BD$ and that $\M(ABCD)$ is regular.
Up to similarity and relabelling, precisely the following three
possibilities occur:
\begin{enumerate}
\item the regular tetrahedron;
\item the disphenoid $T_2$ with
\[
AB=CD=1,\;
AC=AD=BC=BD=\sqrt{\frac{21+4\sqrt6}{45}};
\]
\item a unique nonisosceles class $T_1$, with normalized edges
\[
AB=1,\; CD=\mu_*,\;
AC=AD=BC=BD=\lambda_*,
\]
where
\[
\lambda_*=1.579392054290007551\ldots,\;
\mu_*=2.351892333661864551\ldots.
\]
An exact algebraic specification of $T_1$ is given below.
\end{enumerate}
\end{theorem}

The edge assumption gives the reflections $(1\,2)$ and $(3\,4)$,
which are inherited by the Morley construction. Consequently
\begin{equation}
K=\begin{pmatrix}
1&p&q&q\\p&1&q&q\\q&q&1&s\\q&q&s&1
\end{pmatrix},\;
d=(a,a,b,b)^T,
\label{sym:eq:two}
\end{equation}
where $1<p,q,s<2$ and $a,b>0$.
Choose the regular reference vectors $e_i$ of Appendix~\ref{app:reflection}; the normals are $n_i=\sum_jq_{ij}d_je_j$. Write
\[
U=a(1+p)-2bq,\; V=2aq-b(1+s).
\]
The four normals are
\[
\begin{aligned}
n_1&=(\sqrt2a(1-p),0,U),&
n_2&=(-\sqrt2a(1-p),0,U),\\
n_3&=(0,\sqrt2b(1-s),V),&
n_4&=(0,-\sqrt2b(1-s),V).
\end{aligned}
\]
Comparison of $|n_1-n_2|^2$ and $|n_3-n_4|^2$ gives
\begin{equation}
8a^2=p+2,\; 8b^2=s+2.
\label{sym:eq:ab}
\end{equation}
Their unit lengths and their cross scalar products give
\begin{equation}
U^2=\frac{(p+1)^2(2-p)}4,\;
V^2=\frac{(s+1)^2(2-s)}4,\;
UV=\ph(q).
\label{sym:eq:UV}
\end{equation}

The positive area vector has the form $(\alpha,\alpha,\beta,\beta)$,
so $\alpha U+\beta V=0$. Thus $UV<0$ and $q>\sqrt3$.
Moreover
\[
a(1+p)<\frac3{\sqrt2},\;
2bq>2\sqrt{\frac38}\sqrt3=\frac3{\sqrt2}.
\]
Hence $U<0$ and $V>0$. These signs tell us which square roots to take below.

Set
\begin{equation}
x=\sqrt{\frac{2(2-p)}{p+2}},\;
y=\sqrt{\frac{2(2-s)}{s+2}}.
\label{sym:eq:xy}
\end{equation}
Then
\[
0<x,y<\sqrt{2/3},\;
p=\frac{2(2-x^2)}{2+x^2},\;
s=\frac{2(2-y^2)}{2+y^2},\;
a=\frac1{\sqrt{2+x^2}},\; b=\frac1{\sqrt{2+y^2}}.
\]
Equations \eqref{sym:eq:UV}, with the signs just established, become
\[
2bq=a(p+1)(1+x),\;
2aq=b(s+1)(1+y).
\]
Define
\[
h(z)=\frac{(6-z^2)(1+z)}{(2+z^2)^2}.
\]
We therefore obtain
\begin{equation}
h(x)=h(y)=\gamma,\;
q=\frac{\gamma}{2}\sqrt{(2+x^2)(2+y^2)}.
\label{sym:eq:h}
\end{equation}
The remaining equation $UV=\ph(q)$ is equivalent to
\begin{equation}
q^2-3=\frac{4\gamma xy}{(1+x)(1+y)}.
\label{sym:eq:cross}
\end{equation}
All divisions have positive denominators. Conversely,
\eqref{sym:eq:h} through \eqref{sym:eq:cross} with the indicated ranges reconstruct
the normals and a positive kernel: take face weights proportional
to $(V,V,-U,-U)$. Thus these equations are sufficient as well.

\subsubsection*{The unequal-parameter branch}

Assume $p\ne s$, or equivalently $x\ne y$. Put
\[
u=x+y,\; v=xy,\; 0<v<2/3.
\]
Subtracting $h(x)=h(y)$ and dividing by $x-y$ gives
\begin{equation}
\begin{split}
F(u,v)={}&-6u^3-(6v+4)u^2+(v^2+12v-28)u\\
&+v^3+2v^2-20v+24=0.
\end{split}
\label{sym:eq:F}
\end{equation}
To record the other equation compactly, set
\[
W=4+2u^2-4v+v^2,\;
Z=36-6u^2+12v+v^2,\; T=1+u+v,
\]
and
\[
\begin{split}
C(u,v)={}&6u^4+(6v-4)u^3+(-v^2-24v+20)u^2\\
&+(-v^3-22v^2+36v+24)u+2v^3+4v^2-40v+48.
\end{split}
\]
Here $W=(2+x^2)(2+y^2)$, $Z=(6-x^2)(6-y^2)$,
$T=(1+x)(1+y)$, and
\[
C=(6-x^2)(1+x)(2+y^2)^2
+(6-y^2)(1+y)(2+x^2)^2.
\]
Equation \eqref{sym:eq:cross} is therefore
\begin{equation}
G(u,v):=ZT^2W-12W^2T-8vC(u,v)=0.
\label{sym:eq:G}
\end{equation}

The following identities are checked in \texttt{check\_T1.py}. The resultant is
\begin{equation}
\Res_u(F,G)
=3538944(v-2)^4(v+2)^2(v^2-6)P(v),
\label{sym:eq:resultant}
\end{equation}
where
\begin{equation}
P(v)=v^6-132v^5+4116v^4-33792v^3
+116736v^2-212736v+20736.
\label{sym:eq:P}
\end{equation}
Every factor apart from $P$ is nonzero on $(0,2/3)$.

\begin{lemma}
The polynomial $P$ has exactly one root in $(0,2/3)$.
For this value of $v$, equations $F=G=0$ have exactly one common
value of $u$, and that value is positive.
\end{lemma}

\begin{proof}
For $0<v<2/3$, dropping the negative terms in $P'$ gives
\[
\begin{split}
P'(v)
&<6(2/3)^5+16464(2/3)^3+233472(2/3)-212736\\
&=-\frac{4228928}{81}<0.
\end{split}
\]
Also $P(0)=20736$ and $P(2/3)=-57169664/729$.
This proves the first assertion.

The linear subresultant of $F$ and $G$ is
\[
-6144(v-2)(v+2)\{A(v)u+B(v)\},
\]
where
\[
\begin{split}
A(v)={}&64v^6-2900v^5+31506v^4-98472v^3\\
&+73368v^2+395904v-852672,
\end{split}
\]
\[
\begin{split}
B(v)={}&64v^7-2737v^6+25186v^5-71268v^4\\
&+33144v^3+247008v^2-683136v+614400.
\end{split}
\]
Throughout $(0,2/3)$,
\[
A(v)<64(2/3)^6+31506(2/3)^4+73368(2/3)^2
+395904(2/3)-852672
=-\frac{400876352}{729}<0.
\]
Thus the linear subresultant is nonzero at the root of $P$.
Since the resultant vanishes, the common root is unique and is
\begin{equation}
u=-B(v)/A(v).
\label{sym:eq:u}
\end{equation}
Rational substitution gives
\[
P(0.1031379)>0,\; P(0.1031380)<0.
\]
In particular, $0.10<v<0.11$. On this interval
\[
B(v)>614400-683136(0.11)-71268(0.11)^4
-2737(0.11)^6>0.
\]
Hence $u>0$.
\end{proof}

The unordered pair $\{x,y\}$ is now uniquely determined as the roots
of $z^2-uz+v=0$. Exact interval substitution in \eqref{sym:eq:u} gives
\[
0.67384<u<0.67386,\;
0.2350<\min(x,y)<0.2351,\;
0.4387<\max(x,y)<0.4389.
\]
In particular, the roots are real, positive, distinct, and less than
$\sqrt{2/3}$. Equation \eqref{sym:eq:h} gives $1.84<q<1.85$.
Thus the construction is admissible and proves existence.

This gives an exact description of $T_1$:
take the unique root $v$ of \eqref{sym:eq:P} in $(0,2/3)$, set
$u=-B(v)/A(v)$, let $x<y$ be the roots of $z^2-uz+v$, and recover
$p,s,q$ from \eqref{sym:eq:xy} and \eqref{sym:eq:h}.
Numerically,
\[
\begin{aligned}
v&=0.103137987399601104094765707121\ldots,\\
u&=0.673845625206396229272045732691\ldots,\\
p&=1.892477377432828212380883568822\ldots,\\
s&=1.648726879892260334958056852158\ldots,\\
q&=1.844867969052390261064873920274\ldots .
\end{aligned}
\]
To recover the original edges, put $h_p=\ph(p)$, $h_s=\ph(s)$, and
\[
z^2=\frac{1-h_p}{4(1+h_p)},\;
b_0^2=z^2\frac{1+h_s}{1-h_s}.
\]
The vertices
\[
A=(-1/2,0,0),\; B=(1/2,0,0),\;
C=(0,z,b_0),\; D=(0,z,-b_0)
\]
have the required normal Gram matrix, and
\[
\mu_*=2b_0,\;
\lambda_*=\sqrt{\frac14+z^2+b_0^2}.
\]
Interchanging $x,y$ interchanges the two distinguished opposite edges,
so it gives the same similarity class.


\subsubsection*{The equal-parameter branch}
If $p=s$, the original tetrahedron is a disphenoid. Theorem~\ref{t2:thm:classification} gives precisely the regular tetrahedron and $T_2$. Together with the unequal-parameter branch, this completes the proof of Theorem~\ref{sym:thm:classification}.

\begin{corollary}\label{sym:cor:equivalence}
Conjectures~\ref{conj:four} and~\ref{conj:classification} are equivalent.
\end{corollary}
\begin{proof}
Each of the three classes has four equal cross edges. Conversely, Theorem~\ref{sym:thm:classification} classifies all solutions with that edge property.
\end{proof}


\subsection{Tetrahedra with a nontrivial symmetry}

Combining the last three theorems, we can treat every tetrahedron that has a symmetry.

\begin{corollary}\label{cor:symmetric}
Suppose that $\M(ABCD)$ is regular and that some isometry other than the identity maps the vertex set of $ABCD$ onto itself. Then $ABCD$ has four equal cross edges, and up to similarity it is the regular tetrahedron, $T_1$ or $T_2$. Consequently Conjectures~\ref{conj:four} and~\ref{conj:symmetry} are equivalent, and a counterexample to Conjecture~\ref{conj:four} must have trivial symmetry group.
\end{corollary}

\begin{proof}
The isometry induces a nonidentity permutation $\pi$ of the vertices preserving all six edge lengths. Conversely, every permutation of the vertices preserving the six edge lengths is induced by an isometry, since the edge lengths determine the tetrahedron up to congruence. After relabelling, $\pi$ is one of $(1\,2)$, $(1\,2)(3\,4)$, $(1\,2\,3)$ or $(1\,2\,3\,4)$.

If $\pi=(1\,2)$, the isometry fixes $A_3$, $A_4$ and the midpoint of $A_1A_2$. These three points are not collinear, so the isometry is the reflection in the plane through them. Theorem~\ref{thm:reflection} gives four equal cross edges.

If $\pi=(1\,2)(3\,4)$, then $A_1A_3=A_2A_4$ and $A_1A_4=A_2A_3$. Theorem~\ref{thm:twopairs} applies.

If $\pi=(1\,2\,3)$, then $A_1A_2=A_2A_3=A_3A_1$ and $A_1A_4=A_2A_4=A_3A_4$. The transposition $(1\,2)$ then preserves all edge lengths. It is induced by an isometry, which is a reflection by the first case.

If $\pi=(1\,2\,3\,4)$, then $A_1A_2=A_2A_3=A_3A_4=A_4A_1$ and $A_1A_3=A_2A_4$. All three pairs of opposite edges are equal, and Theorem~\ref{thm:disphenoid} applies.

In every case the tetrahedron has four equal cross edges, and Theorem~\ref{sym:thm:classification} completes the proof.
\end{proof}

\subsection{The arithmetic of the three classes}

The three solutions are quite different arithmetically. For the regular tetrahedron all $q_{ij}$ are equal to the root of $3q^3-9q-2=0$ in $(1,2)$, and $T_2$ is described by square roots (Theorem~\ref{thm:T2}). For $T_1$ the situation is different.

\begin{proposition}\label{prop:T1galois}
The polynomial $P$ in \eqref{sym:eq:P} is irreducible over $\mathbb Q$ and its Galois group is the symmetric group $S_6$. Consequently the number $v$ specifying $T_1$ is not expressible by radicals, and no tetrahedron similar to $T_1$ has vertex coordinates expressible by radicals.
\end{proposition}

\begin{proof}
We use Dedekind's theorem: if a prime $\ell$ divides neither the leading coefficient nor the discriminant of $P$, the degrees of the irreducible factors of $P$ modulo $\ell$ form the cycle type of an element of the Galois group. Modulo $23$, $P$ is irreducible. Modulo $7$ it factors as a quintic times a linear factor, and modulo $31$ as a cubic, a quadratic and a linear factor. The first fact shows that $P$ is irreducible over $\mathbb Q$ and that the Galois group $G$, as a transitive group of degree six, contains a $6$-cycle. The second gives a $5$-cycle fixing a point, so the point stabilizer is transitive on the remaining points; hence $G$ is $2$-transitive and primitive. The cube of an element of type $(3,2,1)$ is a transposition. A primitive permutation group containing a transposition is the full symmetric group. Thus $G=S_6$, which is not solvable, and $v$ is not expressible by radicals.

Suppose a tetrahedron similar to $T_1$ had coordinates expressible by radicals. Its unit normals, and hence $\cos\theta_{ij}=-\ph(q_{ij})$, would be expressible by radicals. Each $q_{ij}$ is a root of $x^3-3x-2\cos\theta_{ij}$, so the numbers $p,s$ of Theorem~\ref{sym:thm:classification} would be expressible by radicals, and so would $x,y$ in \eqref{sym:eq:xy} and $v=xy$, a contradiction.
\end{proof}

The factorizations modulo $7$, $23$ and $31$ are checked in \texttt{check\_galois\_T1.py}. So, unlike $T_2$, the tetrahedron $T_1$ has no expression in radicals, and the description in Theorem~\ref{sym:thm:classification} is the closed form one can hope for.

\section{The case with one pair of equal opposite edges}\label{sec:onepair}

\begin{proposition}
\label{prop:cotangent}
If $AB=CD=\ell$, then
\begin{equation}
24V\ell(\cot\theta_{CD}-\cot\theta_{AB})
=(AC^2-BD^2)(AD^2-BC^2),
\label{eq:cotangent}
\end{equation}
where $\theta_{XY}$ denotes the interior dihedral angle along $XY$. In particular,
\[
\theta_{AB}=\theta_{CD}
\;\Longleftrightarrow\; AC=BD\ \text{or}\ AD=BC.
\]
\end{proposition}

\begin{proof}
Heron's formula for the four faces, using $AB^2=CD^2$, gives
\[
8(S_A^2+S_B^2-S_C^2-S_D^2)
=(AC^2-BD^2)(AD^2-BC^2).
\]
For a direct check, put $a=\ell^2$, $x=AC^2$, $y=AD^2$, $z=BC^2$, $t=BD^2$. In the corresponding sum of the four expressions for $16S_i^2$, only $2(zt+xy-yt-xz)=2(x-t)(y-z)$ remains.

The equilibrium of the area normals gives
\[
S_A^2+S_B^2-S_C^2-S_D^2
=2S_AS_B\cos\theta_{CD}-2S_CS_D\cos\theta_{AB}.
\]
By \eqref{eq:edge}, both $S_AS_B\sin\theta_{CD}$ and $S_CS_D\sin\theta_{AB}$ equal $3V\ell/2$. Substitution proves \eqref{eq:cotangent}. The last assertion follows because cotangent is injective on $(0,\pi)$.
\end{proof}

Together with Theorem~\ref{thm:twopairs}, the proposition shows that Conjecture~\ref{conj:onepair} would follow from
\begin{equation}
\M(ABCD)\text{ regular},\; AB=CD
\;\Longrightarrow\; \theta_{AB}=\theta_{CD}.
\label{eq:missing}
\end{equation}
Conversely, suppose that $AB=CD$ and that the tetrahedron has four equal cross edges with respect to some pair of opposite edges. Then it has a second pair of equal opposite edges, and Proposition~\ref{prop:cotangent} gives $\theta_{AB}=\theta_{CD}$. Hence, for tetrahedra with $AB=CD$, Conjecture~\ref{conj:onepair} is equivalent to \eqref{eq:missing}; the implication itself is not proved here. With $A,B,C,D=A_1,A_2,A_3,A_4$, it asks for $q_{12}=q_{34}$.

\section{Near the regular tetrahedron}\label{sec:nearregular}

\begin{theorem}
\label{thm:localregular}
There is a neighborhood of the regular shape in the space of tetrahedra modulo similarities in which a regular Morley tetrahedron implies that the original tetrahedron is regular.
\end{theorem}

\begin{proof}
Let
\[
R=\begin{pmatrix}1&1&1\\1&-1&-1\\-1&1&-1\\-1&-1&1\end{pmatrix},
\]
with rows $r_i$. Parametrize centered tetrahedra by the rows of $X=RG^{1/2}$, where $G$ is symmetric positive definite and $\tr G=3$. The regular shape is $G=I_3$. Put
\[
\mu_i=(r_iG^{-1}r_i^T)^{1/2},\;
H_{ij}=\frac{r_iG^{-1}r_j^T}{\mu_i\mu_j},
\]
and obtain $K$ from $H$ as in Section~\ref{sec:gram}. The $\mu_i$ are proportional to the face areas. The barycentric coordinate matrix of the Morley vertices is
\[
P(G)=\diag\bigl((K\mu)_i^{-1}\bigr)K\diag(\mu_i).
\]
Consequently, if $M=P(G)X$ and $Y=R^TM/4$, the normalized shape of the Morley tetrahedron is
\begin{equation}
\Phi(G)=\frac{3YY^T}{\tr(YY^T)}.
\label{eq:shapemap}
\end{equation}
The map is real analytic near $I_3$, and the Morley tetrahedron is regular exactly when $\Phi(G)=I_3$.

Let
\[
q=2\cos\left(\frac13\arccos\frac13\right),\;
3q^3-9q-2=0,
\]
and put $\lambda=1+3q$, $h=(1-q)/\lambda$. At $G=I_3$ we have
\[
H_0=(4I_4-\one\one^T)/3,\;
K_0=(1-q)I_4+q\one\one^T,\; Y_0=hI_3.
\]
Let $E$ be a symmetric matrix with trace zero. Write it as the sum of a diagonal matrix and a matrix whose diagonal entries are zero. Differentiation gives
\begin{equation}
D\Phi_{I_3}(E)=\kappa_2 E_{\mathrm{diag}}+\kappa_3 E_{\mathrm{off}},
\label{eq:differential}
\end{equation}
where
\begin{equation}
\kappa_2=\frac{8-9(q-1)^2}{9(q-1)^2(q+1)},
\;
\kappa_3=-\frac{8-9(q-1)^2}{27(q-1)^2(1+3q)}.
\label{eq:kappas}
\end{equation}

For completeness, the intermediate differentials are as follows. Set $a_i=r_iEr_i^T$, $v_i=-a_i/6$, and $\eta=-2/[3(q^2-1)]$. Then
\[
\delta\mu_i/\sqrt3=v_i,\;
\delta H_{ij}=-\frac{r_iEr_j^T}{3}+(H_0)_{ij}\frac{a_i+a_j}{6},
\]
and $\delta K_{ij}=\eta\delta H_{ij}$ off the diagonal, with zero diagonal. Also,
\[
\delta P=\frac{\delta K+K_0\diag v}{\lambda}
-\frac{\diag(\delta K\one+K_0v)K_0}{\lambda^2},
\;
\delta Y=\frac14R^T\delta P R+\frac h2E.
\]
Substitution into \eqref{eq:shapemap}, evaluated on $\diag(1,-1,0)$ and on the symmetric matrix with entries $E_{23}=E_{32}=1$ and all others zero, gives respectively
\[
1-\frac{16}{9(q-1)^2(q+1)},\;
\frac{q+3}{3(1+3q)}-\frac{32}{27(q-1)^2(1+3q)}.
\]
Using $3q^3-9q-2=0$ reduces these to \eqref{eq:kappas}; coordinate permutations give the remaining basis directions and establish \eqref{eq:differential}.

The cubic is strictly increasing on $(1,2)$ and changes sign between $9/5$ and $19/10$. Hence $8-9(q-1)^2>0$, so $\kappa_2>0$ and $\kappa_3<0$. The two tangent subspaces have dimensions $2$ and $3$, respectively, and the differential is invertible on their direct sum. The inverse function theorem now implies that $\Phi(G)=I_3$ has only the solution $G=I_3$ in a neighborhood of $I_3$.
\end{proof}

\section{Summary of the proved cases}
\label{sec:summary}

The table lists what is known about a tetrahedron $ABCD$ with a regular Morley tetrahedron under various additional hypotheses.

\begin{center}\small
\begin{tabular}{p{.43\textwidth}p{.47\textwidth}}\toprule
Additional hypothesis & Conclusion\\\midrule
$ABCD$ isosceles & regular or $T_2$ (Theorems~\ref{thm:disphenoid}, \ref{t2:thm:classification})\\
$t_1=t_2=t_3=t_4$ & isosceles, hence regular or $T_2$ (Theorem~\ref{thm:equaldistances})\\
two pairs of equal opposite edges & four equal cross edges (Theorem~\ref{thm:twopairs})\\
one reflection symmetry & a second reflection (Theorem~\ref{thm:reflection})\\
four equal cross edges & regular, $T_1$ or $T_2$ (Theorem~\ref{sym:thm:classification})\\
any nontrivial symmetry & regular, $T_1$ or $T_2$ (Corollary~\ref{cor:symmetric})\\
shape close to the regular one & regular (Theorem~\ref{thm:localregular})\\
shape close to $T_1$ or $T_2$ & equal to it (Section~\ref{sec:intervals})\\
$AB=CD$ only & open; equivalent to $\theta_{AB}=\theta_{CD}$ (Proposition~\ref{prop:cotangent})\\
none & open (Conjecture~\ref{conj:four})\\\bottomrule
\end{tabular}
\end{center}

What remains open is the case of a tetrahedron with trivial symmetry group that is not close to one of the three known solutions.

\section{Possible approaches}
\label{tet:directions}

By Corollaries~\ref{sym:cor:equivalence} and~\ref{cor:symmetric}, everything comes down to one step: to show that a regular Morley tetrahedron forces some symmetry of the original tetrahedron. We describe several ways one might attempt this.

\subsection{Forcing a first reflection}

By Proposition~\ref{prop:defect}, a reflection symmetry corresponds to equality $\varepsilon_{ij}=0$ in a nonnegative quantity. One could try to combine the six defects with the positive kernel condition for $H$ and the identity \eqref{eq:regularity}, and to show that at least one defect vanishes; Theorem~\ref{thm:reflection} would then give the second reflection. It would suffice to find an identity or inequality that rules out $\varepsilon_{ij}>0$ for all six pairs simultaneously. The bounds of Proposition~\ref{prop:distancebounds} restrict the variables, but they do not come close to this.

\subsection{One pair of equal opposite edges}

The case $AB=CD$ is a natural test. Here one has to prove \eqref{eq:missing}, that is, $q_{12}=q_{34}$; Theorem~\ref{thm:twopairs} then finishes the argument. One can try to combine the edge condition \eqref{eq:edge} with some of the entries of \eqref{eq:regularity}. Any such elimination has to keep track of the positivity of the face areas, since polynomial relations alone have roots that do not come from tetrahedra.

\subsection{The case of acute dihedral angles}

The exceptional examples both have obtuse dihedral angles: $T_1$ has one and $T_2$ has two. It is conceivable that a tetrahedron with six acute dihedral angles and a regular Morley tetrahedron must be regular. In the matrix formulation acuteness means $\ph(q_{ij})<0$ for all $i\ne j$, and this sign condition might allow estimates that are not available in general.

\subsection{A global computation}

The system $F=0$ of \eqref{eq:tenequations} consists of ten polynomial equations of degree at most four in the ten unknowns $q_{ij}$ and $d_i$. It is invariant under the action of $S_4$ and under $d\mapsto-d$, and its admissible solutions lie in a bounded region: $1<q_{ij}<2$, $0<d_i\le1/\sqrt2$ by Proposition~\ref{prop:defect}, and $\max_id_i/\min_id_i<2/\sqrt3$ by Proposition~\ref{prop:distancebounds}. This suggests two computational approaches.

First, one can cover the closure of this region by boxes and exclude zeros by rational interval arithmetic, as in Section~\ref{sec:intervals}, except in small neighbourhoods of the three known solutions. The difficulty is the boundary. When some $q_{ij}$ tends to $1$ or $2$ the tetrahedron degenerates, and one needs to know which variables can approach their limits and whether the equations can hold along such a sequence. This requires estimates that remain valid near degeneration.

Second, homotopy continuation or monodromy methods could be used to compute all isolated complex solutions of $F=0$, keeping the real solutions that satisfy the conditions of Proposition~\ref{prop:regularity}. The B\'ezout bound $4^{10}$ is large, but the symmetries help, and by Corollary~\ref{cor:symmetric} any new solution would appear in $24$ labelled copies. Such a computation would give strong evidence; a proof would require certified path tracking or a separate argument for completeness.

\subsection{The sector construction}

For the generalized sector construction our methods do not apply directly. The trisection case relies on the cubic relation $\ph(q)=(3q-q^3)/2$ between $K$ and $H$. For other sector parameters one would first have to find the corresponding distance ratios and see which parts of the argument survive. We have not studied Conjecture~\ref{conj:psector} beyond the trisection case.

\section{Open problems}
\label{sec:openproblems}

We end with a list of questions.
\begin{enumerate}[label=\arabic*.]
\item Prove or disprove Conjecture~\ref{conj:four}, or one of its equivalent forms: the classification Conjecture~\ref{conj:classification}, the symmetry Conjecture~\ref{conj:symmetry}, or the matrix form Conjecture~\ref{conj:positivity}.
\item Prove Conjecture~\ref{conj:onepair}, that is, the implication \eqref{eq:missing}.
\item Determine all tetrahedra whose Morley tetrahedron is isosceles. By Theorem~\ref{thm:curve} they include a curve of nonisosceles tetrahedra through $T_1$. Does every such tetrahedron have a symmetry?
\item Must a tetrahedron with six acute dihedral angles and regular Morley tetrahedron be regular?
\item Find the analogue of the Gram formulation for the sector construction of \cite[Theorem 6]{TranActa}, and decide Conjecture~\ref{conj:psector}.
\item Is there a geometric reason why $T_2$ is expressible by square roots while $T_1$ is defined by a sextic with Galois group $S_6$?
\item Study the same problems for simplices of higher dimension.
\end{enumerate}

\begin{appendices}

\section{A polynomial inequality for two opposite pairs}
\label{app:bernstein}

This appendix proves the incompatibility used in Theorem~\ref{thm:twopairs}. Set
\[
x=P,\; t=R/P,\;
z=\frac{wP(P+3)}{R(R+3)}.
\]
By \eqref{eq:wbound}, all three variables lie in $(0,1)$. Write
\[
u=x(1+t),\; v=tx^2,\; D_0=x+3,
\; W=zt(tx+3),\; w=W/D_0.
\]
Define the polynomials
\begin{align*}
N={}&2u(2u+3)D_0^2-6uWD_0-v(W^2+4WD_0+11D_0^2),\\
C_0={}&v(u+5)D_0+v(u+1)W+2Wu(u+3),\\
E={}&\frac{NC_0-6(u+1)(v+2u+4)v(D_0-W)(D_0+W)^2}{x^2},\\
B_0={}&(u^2+3u-6)D_0^2-6WD_0+2v(W^2+WD_0-D_0^2).
\end{align*}
The numerator defining $E$ is divisible by $x^2$. Equation \eqref{eq:Ezero} becomes $E=0$, while \eqref{eq:Bpositive} becomes $B_0>0$.

\begin{lemma}
\label{lem:bernstein}
On $(0,1)^3$, $B_0\ge0$ implies $E<0$.
\end{lemma}

\begin{proof}
Let $b_{n,i}(X)=\binom niX^i(1-X)^{n-i}$, and define
\[
Q(x,t,z)=\sum_{i=0}^2\sum_{j=0}^4\sum_{k=0}^1
q_{ijk}b_{2,i}(x)b_{4,j}(t)b_{1,k}(z).
\]
The nonzero coefficients are the following positive rational numbers; all other coefficients are zero.

\begin{center}
\begin{tabular}{cc}
\toprule
$(i,j,k)$ & $q_{ijk}$\\
\midrule
$(0,1,1)$ & $30902484965806161/7009635898549648$\\
$(0,2,1)$ & $16399650927464661/7009635898549648$\\
$(0,3,1)$ & $129949208779311819/14019271797099296$\\
$(0,4,0)$ & $48/7$\\
$(0,4,1)$ & $24/7$\\
$(1,1,1)$ & $2483086725481191/7009635898549648$\\
$(1,3,1)$ & $382341243899593387/56077087188397184$\\
$(1,4,0)$ & $869/56$\\
$(2,0,1)$ & $386384/44499$\\
$(2,1,1)$ & $512549/44499$\\
$(2,2,1)$ & $2202979/38142$\\
$(2,4,0)$ & $73/6$\\
\bottomrule
\end{tabular}
\end{center}

Put $R_0=-E-B_0Q$. In the tensor Bernstein basis of degree $(6,9,3)$, all $280$ coefficients of $R_0$ are nonnegative, and $252$ are strictly positive. These signs follow by substituting the displayed coefficients into the following formula. Namely, if $R_0=\sum_\nu a_\nu X^\nu$, the coefficient at multi-index $\iota$ is
\[
r_\iota=\sum_{\nu\le\iota}a_\nu
\prod_{j=1}^3\frac{\binom{\iota_j}{\nu_j}}{\binom{n_j}{\nu_j}},
\; n=(6,9,3).
\]
Thus each sign follows from a finite sum of the rational numbers given above.

Every Bernstein basis function is positive in the open cube. Hence $R_0>0$ there, while $Q\ge0$. The identity $-E=B_0Q+R_0$ proves the lemma and completes the polynomial step of Theorem~\ref{thm:twopairs}.
\end{proof}

\section{Polynomial elimination for one reflection}
\label{app:reflection}

We use the notation from the proof of Theorem~\ref{thm:reflection}. Put
\[
e_1=(\sqrt2,0,1),\; e_2=(-\sqrt2,0,1),\;
e_3=(0,\sqrt2,-1),\; e_4=(0,-\sqrt2,-1).
\]
Choose $M_i=Le_i/(2\sqrt2)$. The normals satisfy $n_i=\sum_jq_{ij}d_je_j$, giving
\begin{align*}
n_1&=(\sqrt2a(1-p),\sqrt2(bq-cr),a(1+p)-bq-cr),\\
n_3&=(0,\sqrt2(b-cs),2aq-b-cs),\\
n_4&=(0,\sqrt2(bs-c),2ar-bs-c).
\end{align*}
The difference $n_1-n_2$ gives $a^2=(p+2)/8$.

Assume $b\ne c$, and relabel so that $b>c$. Define
\[
R=\frac{b+c}{a},\; U=\frac{b-c}{a},\; X=R^2,\; Y=U^2.
\]
Thus $X>Y>0$. The vectors $n_3+n_4$ and $n_3-n_4$ lie in the $yz$ plane, are perpendicular, and have squared lengths $(s+1)^2(2-s)$ and $(s-1)^2(s+2)$. Comparison of their coordinates yields a signed parameter $T$ such that
\[
T^2=\frac{2-s}{2(s+2)},\;
s=\frac{2(1-2T^2)}{1+2T^2},\; 0<|T|<1/\sqrt6,
\]
and
\[
q+r=\alpha R,\; q-r=\beta U,\;
\alpha=\frac{(s+1)(1+2T)}2,\;
\beta=\frac{(s-1)(T^{-1}-1)}2.
\]
The squared lengths also give
\begin{equation}
(p+2)\left(\frac{X}{s+2}+\frac{Y}{2-s}\right)=4.
\label{eq:reflectionlength}
\end{equation}

The sign of $T$ is forced by the positive kernel. Set
\[
D_* = p+1-\frac{\alpha^2X}{s+1}+\frac{\beta^2Y}{s-1}.
\]
Block elimination gives $\det K=(1-p)(1-s^2)D_*$, so $D_*<0$ by Lemma~\ref{lem:negative}. Rescale the positive area vector so that
\[
KS=(1,1,a/b,a/c)^T.
\]
Solving this system and using \eqref{eq:reflectionlength}, its first coordinate is
\[
S_1=S_2=-\frac{8T(s+2)}{(p+2)(X-Y)D_*}.
\]
Since $S_1>0$, we must have $T>0$. The physical parameter interval is therefore $(0,1/\sqrt6)$.

\subsection{The necessary rational system}

Put $\gamma=(s+1)T$ and $\delta=(s-1)/(2T)$. Equation \eqref{eq:reflectionlength} is the case $i=0$ of the following system. The cases $i=1,2$ are obtained from the equations $n_1\cdot(n_3+n_4)=\ph(q)+\ph(r)$ and $n_1\cdot(n_3-n_4)=\ph(q)-\ph(r)$, using $|n_1|=1$:
\begin{equation}
A_iX+B_iY+C_i=0\;(i=0,1,2),
\label{eq:reflectionlinear}
\end{equation}
where
\[
A_0=(p+2)/(s+2),\; B_0=(p+2)/(2-s),\; C_0=-4,
\]
\begin{align*}
A_1&=-(p+2)\gamma\alpha+\alpha^3,\\
B_1&=(p+2)((\alpha+\beta)(s+1)-\gamma\beta)+3\alpha\beta^2,\\
C_1&=2\gamma(p+2)(p+1)-12\alpha,\\
A_2&=(p+2)(-(\alpha+\beta)(s-1)-\delta\alpha)+3\beta\alpha^2,\\
B_2&=-(p+2)\delta\beta+\beta^3,\\
C_2&=2\delta(p+2)(p+1)-12\beta.
\end{align*}
The last equation uses $U\ne0$. The unit length of $n_1$ gives
\begin{equation}
(p+2)\left[2(p-1)^2+\frac{(\alpha+\beta)^2XY}{2}
+\left(p+1-\frac{\alpha X+\beta Y}{2}\right)^2\right]=8.
\label{eq:reflectionunit}
\end{equation}

For $j=1,2$, set
\[
D_j=A_0B_j-A_jB_0,\;
X_j=B_0C_j-B_jC_0,\;
Y_j=A_jC_0-A_0C_j.
\]
Every solution of \eqref{eq:reflectionlinear} satisfies $X_j=D_jX$ and $Y_j=D_jY$, including when $D_j=0$. It therefore satisfies $F=G_1=G_2=0$, where
\[
F=D_1C_2+A_2X_1+B_2Y_1
\]
and
\begin{align*}
G_j={}&(p+2)\bigl[8(p-1)^2D_j^2+2(\alpha+\beta)^2X_jY_j\\
&\hspace{2em}+(2(p+1)D_j-\alpha X_j-\beta Y_j)^2\bigr]-32D_j^2.
\end{align*}
No division by $D_j$ has been made, so this necessary system loses no singular case.

\subsection{The two resultants}

We describe the polynomial calculation explicitly. For a polynomial in $p$ with coefficients in $\mathbb Q(T)$, let $\pp$ mean the following operation: clear the denominators, divide the coefficients by their polynomial greatest common divisor in $\mathbb Q[T]$, and multiply by a rational constant so that the resulting integer coefficients have greatest common divisor $1$. The result is unique up to sign. Put
\[
\widetilde F=\pp(F),\; \widetilde G_j=\pp(G_j)\;(j=1,2).
\]
To divide without introducing new denominators, start with $V=\widetilde G_j$ and, while $\deg_p V\ge3$, replace $V$ by
\[
\operatorname{lc}_p(\widetilde F)V
-\operatorname{lc}_p(V)p^{\deg_p V-3}\widetilde F.
\]
Let $H_j=\pp(V)$ at the end of this procedure. Substitution of the preceding formulas gives $\deg_p\widetilde F=3$ and $\deg_pH_j=2$.

The denominators, polynomial factors removed by these operations, and the leading coefficient of $\widetilde F$ have only the following possible nonconstant factors:
\[
T,\; 1+2T^2,\; 2T^2-3,\; 6T^2-1,\; 6T^2-4T-3.
\]
None vanishes for $0<T<1/\sqrt6$. Consequently every solution under consideration gives a common zero of $\widetilde F,H_1,H_2$.

For a cubic $a_3p^3+a_2p^2+a_1p+a_0$ and a quadratic $b_2p^2+b_1p+b_0$, define
\[
\begin{aligned}
\mathcal A&=a_1b_2^2-a_3b_0b_2-a_2b_1b_2+a_3b_1^2,\\
\mathcal B&=a_0b_2^2-a_2b_0b_2+a_3b_1b_0.
\end{aligned}
\]
Their resultant is the polynomial quotient
\[
\frac{b_0\mathcal A^2-b_1\mathcal A\mathcal B+b_2\mathcal B^2}{b_2^2}.
\]
Use this formula to define $R_j(T)=\Res_p(\widetilde F,H_j)$ for $j=1,2$. Thus both $R_j$ are determined by the rational functions already given. Up to a nonzero constant, the factorization of $R_1$ is
\begin{align*}
R_1={}&(T-1)^8(2T+1)^{12}(2T^2+1)^{12}(2T^2-3)^2(6T^2-1)^2\\
&\cdot(6T^2-4T-3)^{10}(2T^2+8T-1)^4\mathcal P(T)\mathcal E(T)^2.
\end{align*}
Here $\mathcal P$ and $\mathcal E$ are specified by their coefficients in ascending order. The coefficients of $\mathcal P$, from $T^0$ through $T^{18}$, are
\[
\begin{gathered}
(-243,2349,-6561,5229,140586,-485388,-1094108,\\
3265740,5130696,-10136384,-10261392,13062960,\\
8752864,-7766208,-4498752,334656,839808,601344,124416).
\end{gathered}
\]
The coefficients of $\mathcal E$, from $T^0$ through $T^{19}$, are
\[
\begin{gathered}
(27,-2430,8190,60228,26352,-164448,-3342112,-8861120,\\
22756256,54892992,-100341696,-121096320,264964864,\\
20152832,-294750720,214062080,-35962112,-18583040,\\
4288000,1664000).
\end{gathered}
\]
The Euclidean algorithm gives $\gcd(\mathcal E,R_2)=1$. All the other factors of $R_1$, apart from $2T^2+8T-1$ and $\mathcal P$, have no zero in $(0,1/\sqrt6)$. Thus a solution must satisfy one of these last two equations.

At the positive root of $2T^2+8T-1$, the polynomial $\widetilde F$ is proportional to
\[
(p-1)(p-9/7)(p-17/3).
\]
Only $p=9/7$ lies in $(1,2)$. It fails $G_1=0$: substitution in the rational expression for $G_1$ gives a numerator with a nonzero linear remainder modulo $2T^2+8T-1$, and a denominator that does not vanish. Since the positive root has degree two over $\mathbb Q$, this excludes the first possibility.

For the second possibility, form the Sturm sequence beginning with $\mathcal P,\mathcal P'$ and continuing by taking the negative remainder at each Euclidean division. Multiplication of a term by a positive constant does not change its signs. The variation counts at the following rational points are
\begin{center}
\begin{tabular}{crrrrrr}
\toprule
$T$ & $0$ & $0.143$ & $0.144$ & $0.295$ & $0.296$ & $0.409$\\
\midrule
Variations & $9$ & $9$ & $8$ & $8$ & $7$ & $7$\\
\bottomrule
\end{tabular}
\end{center}
Since $0.296<1/\sqrt6<0.409$, precisely two roots remain. For the pair $\widetilde F,H_1$, the Euclidean algorithm also gives
\[
\gcd(\mathcal P,b_2)=\gcd(\mathcal P,\mathcal A)=1.
\]
At either root, therefore, the last linear remainder yields $p=-\mathcal B/\mathcal A$. Substitution into $D_1$ gives a nonzero value at both roots, so $X=X_1/D_1$ and $Y=Y_1/D_1$.

It remains to apply positive face areas. Write
\[
h=\ph(p),\; u=\ph(q),\; v=\ph(r),\; w=\ph(s),
\; k=(1+h)/2.
\]
The vector $N=(n_1+n_2)/2$ has squared length $k$, and $N\cdot n_3=u$, $N\cdot n_4=v$. The equilibrium equation
\[
2S_1N+S_3n_3+S_4n_4=0
\]
gives the necessary inequality
\begin{equation}
\mathcal C:=uv-wk=\frac{(1-w^2)S_3S_4}{4S_1^2}>0.
\label{eq:reflectionpositive}
\end{equation}
Set
\[
Z=\frac{\alpha^2X-\beta^2Y}{4}=qr,\;
W=\frac{\alpha^2X+\beta^2Y}{2}=q^2+r^2.
\]
Then
\[
4\mathcal C=Z(9-3W+Z^2)-\ph(s)(2+3p-p^3),
\]
a rational expression in $T,p,X,Y$. Start with the two intervals for $T$ in the preceding table and bisect, retaining the unique root counted by the Sturm sequence, until the width is less than $10^{-70}$. Substitution in the formulas for $p,X,Y$ and $4\mathcal C$, using rational intervals, gives
\begin{center}
\begin{tabular}{ccc}
\toprule
Root interval for $T$ & Interval for $p$ & Interval for $4\mathcal C$\\
\midrule
$(0.143,0.144)$ & $(1.69702,1.69703)$ & $(-0.704,-0.703)$\\
$(0.295,0.296)$ & $(1.25214,1.25215)$ & $(-0.168,-0.167)$\\
\bottomrule
\end{tabular}
\end{center}
Both possibilities contradict \eqref{eq:reflectionpositive}. Thus $b\ne c$ is impossible, completing the algebraic step of Theorem~\ref{thm:reflection}.
\section{Verification of the computations}\label{sec:reproducibility}

The computer-assisted steps of the paper are checked by the following Python scripts, which are available as ancillary files with the arXiv version. They require only the Python standard library, and all proof-critical comparisons use exact rational arithmetic. Decimal calculations in \texttt{check\_T1.py} are used only to print approximations. Each script names the statements it checks, in the order in which they appear in the text. Polynomial identities in one or two variables are checked by evaluation at sufficiently many rational points, using an a priori bound for the degrees.

\begin{center}\small
\begin{tabular}{p{.33\textwidth}p{.59\textwidth}}\toprule
File&Content\\\midrule
\texttt{check\_section6.py}&the bounds \eqref{eq:contractionbounds} and \eqref{eq:intervalgeometry} for the three boxes, the localization \eqref{eq:sharpbox}, the edge bounds and the bound for $\partial f/\partial b$ (Section~\ref{sec:intervals})\\
\texttt{check\_T2.py}&the exact identities in the proof of Theorem~\ref{thm:T2} and the resultant \eqref{t2:eq:resultant}\\
\texttt{check\_T1.py}&the identities, the resultant \eqref{sym:eq:resultant}, the linear subresultant, the bounds of the lemma, the admissibility intervals and the decimal values in Theorem~\ref{sym:thm:classification}\\
\texttt{check\_galois\_T1.py}&the factorizations used in Proposition~\ref{prop:T1galois}\\
\texttt{check\_appendixA.py}&the $280$ Bernstein coefficients of Appendix~\ref{app:bernstein}\\
\texttt{check\_appendixB.py}&the elimination, the factorization of $R_1$, the Sturm counts and the final intervals of Appendix~\ref{app:reflection}\\
\texttt{morley\_tools.py}&polynomial arithmetic, rational intervals, determinants, Sturm sequences, factorization modulo a prime\\\bottomrule
\end{tabular}
\end{center}

The three contraction constants in \eqref{eq:contractionbounds} are less than $.061689$, $.059243$ and $.065397$. Among the $280$ Bernstein coefficients in Appendix~\ref{app:bernstein}, $28$ vanish and $252$ are positive.

\end{appendices}

\section*{Acknowledgements}
\phantomsection
\addcontentsline{toc}{section}{Acknowledgements}

The author thanks Tran Manh Dung and Nguyen Xuan Tho for their help with the computations. The AI assistants ChatGPT and Claude were used for exploratory computations, for checking algebra and in preparing the manuscript; the author takes full responsibility for the content.

\phantomsection
\addcontentsline{toc}{section}{References}
\begin{thebibliography}{98}

\bibitem{BrandtsKrizek}
J.~Brandts and M.~K\v{r}\'i\v{z}ek, \emph{Duality of isosceles tetrahedra}, J. Geom. \textbf{110} (2019), Article 49, 11 pp.

\bibitem{Braude}
E.~J. Braude, \emph{Generalizing the Morley trisector and various theorems with realizability computations}, arXiv:1603.03463, 2016.

\bibitem{Child}
J.~M. Child, \emph{Proof of Morley's theorem}, Math. Gaz. \textbf{11} (1923), 171.

\bibitem{CoxeterGreitzer}
H.~S.~M. Coxeter and S.~L. Greitzer, \emph{Geometry Revisited}, New Mathematical Library, Vol.~19, Mathematical Association of America, Washington, DC, 1967.

\bibitem{DergiadesTran}
N.~Dergiades and Q.~H. Tran, \emph{On some extensions of Morley's trisector theorem}, J. Geom. Graph. \textbf{24} (2020), 197--205.

\bibitem{GrinbergOrhon}
E.~L. Grinberg and M.~Orhon, \emph{Morley trisectors and the law of sines with reflections}, Amer. Math. Monthly \textbf{128} (2021), 163--167.

\bibitem{HajjaWalker}
M.~Hajja and P.~Walker, \emph{Equifacial tetrahedra}, Int. J. Math. Educ. Sci. Technol. \textbf{32} (2001), 501--508.

\bibitem{Hashimoto}
Y.~Hashimoto, \emph{A short proof of Morley's theorem}, Elem. Math. \textbf{62} (2007), 121.

\bibitem{Klain}
D.~A. Klain, \emph{Tetrahedra with congruent facet pairs}, Math. Intelligencer \textbf{45} (2023), 251--255.

\bibitem{Krawczyk}
R.~Krawczyk, \emph{Newton-Algorithmen zur Bestimmung von Nullstellen mit Fehlerschranken}, Computing \textbf{4} (1969), 187--201.

\bibitem{Leech}
J.~Leech, \emph{Some properties of the isosceles tetrahedron}, Math. Gaz. \textbf{34} (1950), 269--271.

\bibitem{Morley}
F.~Morley, \emph{On the metric geometry of the plane $n$-line}, Trans. Amer. Math. Soc. \textbf{1} (1900), 97--115.

\bibitem{Neumaier}
A.~Neumaier, \emph{Interval Methods for Systems of Equations}, Encyclopedia of Mathematics and its Applications, Vol.~37, Cambridge University Press, Cambridge, 1990.

\bibitem{OakleyBaker}
C.~O. Oakley and J.~C. Baker, \emph{The Morley trisector theorem}, Amer. Math. Monthly \textbf{85} (1978), 737--745.

\bibitem{Rump}
S.~M. Rump, \emph{Verification methods: rigorous results using floating point arithmetic}, Acta Numer. \textbf{19} (2010), 287--449.

\bibitem{Schneider}
R.~Schneider, \emph{Convex Bodies: The Brunn--Minkowski Theory}, 2nd expanded ed., Encyclopedia of Mathematics and its Applications, Vol.~151, Cambridge University Press, Cambridge, 2014.

\bibitem{SvrtanVeljan}
D.~Svrtan and D.~Veljan, \emph{Side lengths of Morley triangles and tetrahedra}, Forum Geom. \textbf{17} (2017), 123--142.

\bibitem{TranActa}
Q.~H. Tran, \emph{Morley's trisector theorem for isosceles tetrahedron}, Acta Math. Hungar. \textbf{165} (2021), 308--315. \href{https://doi.org/10.1007/s10474-021-01184-0}{doi:10.1007/s10474-021-01184-0}.

\bibitem{TranSynthetic}
Q.~H. Tran, \emph{A synthetic proof of the Morley trisector theorem using congruent and similar triangles}, Elem. Math. \textbf{80} (2025), 70--72.

\bibitem{TranVolume}
Q.~H. Tran, \emph{A Generalized Volume Formula for Tetrahedra with
Congruent Facet Pairs}, Math. Intelligencer \textbf{47} (2025),
136--139. \href{https://doi.org/10.1007/s00283-024-10375-w}
{doi:10.1007/s00283-024-10375-w}.

\end{thebibliography}

\begin{flushright}
Quang Hung Tran,\\
High School for Gifted Students,\\
Vietnam National University, Hanoi, Vietnam,\\
Email: \href{mailto:tranquanghung@hus.edu.vn}{tranquanghung@hus.edu.vn}\\
ORCID: \href{https://orcid.org/0000-0003-2468-4972}{0000-0003-2468-4972}
\end{flushright}

\end{document}
